\chapter{Reading Lean, and a map of the book}\label{ch:01}
\begingroup\setlength{\emergencystretch}{3em}

\providecommand{\chonevac}{\mathbf{1}}
\providecommand{\choneV}{\mathcal{V}}

This book documents the release of the LeanVOA library named in the front matter. Every Lean quotation in this chapter is quoted from the files of that
release, and every one is machine-checked by \texttt{scripts/textbook/sync.py} to occur verbatim
(up to whitespace and comments) in the source. This first chapter contains almost no vertex
algebra. Its job is to prepare the reader for the rest of the book: to say what the book is for
and how to use it, to collect in one place every Lean and Mathlib idiom that the later chapters use
(so that their fidelity checks can simply point back here, Section~\ref{ch01:sec:glossary}), to explain exactly what a computer
check does and does not guarantee, to fix the conventions that recur in every chapter, and to give a
map of the book and of where each basic notion is defined. The mathematics proper begins in
Chapter~\ref{ch:L01}.

%======================================================================
\section{What this book is for, and how to use it}\label{ch01:sec:purpose}
%======================================================================

\subsection{The referee's question}

The LeanVOA library is a body of roughly ten thousand Lean files (the per-chapter counts at the book's commit are tabulated in
Appendix~\ref{ch:provenance})
which formalise a large part of the theory of vertex operator algebras and, as its headline result,
construct the Moonshine module $\Vnat$ of Frenkel, Lepowsky and Meurman. The Lean system has checked
every proof in it. That is a strong guarantee, but it is a guarantee about \emph{proofs}, not about
\emph{statements}. A Lean theorem can be completely verified and still not say what its authors
meant: a hypothesis may be stronger than intended, a conclusion weaker, a definition may be vacuous
(satisfied by nothing, or by a trivial object), or a convention may have silently changed between
two files. No computer can detect such a mismatch, because the ``intended meaning'' lives in the
heads of the authors and in the literature, not in the code.

This book exists to close that gap. It is written for a \emph{human referee}: a reader who wants
to take a theorem of the library, read it as mathematics, read the Lean declaration that is
printed next to it, and convince herself that the Lean statement is \emph{exactly} the intended
theorem --- no stronger hidden hypotheses, no weaker conclusion, no vacuous definitions, no changed
conventions. The book is about statements and their meaning. It is not about how Lean checks
proofs; the proofs are explained on paper, as a mathematician (and, where it helps, a physicist)
would explain them.

\subsection{How each result is presented}

Every definition and theorem in the later chapters is presented in the same order:
\begin{enumerate}
\item the \textbf{mathematical statement}, in textbook form, with every quantifier explicit;
\item the \textbf{Lean declaration}, copied verbatim from the source file, with its full name and
  its file (the line ``File \leanpath{VOA/...}'' above or below the code);
\item a \textbf{fidelity check} (a grey box) which decodes the Lean notation, lists every
  hypothesis and says why it is the intended one, traces every defined term back to its
  definition, says whether the Lean form is the same as, stronger than or weaker than the textbook
  form, and asks how the statement could have been wrong or vacuous;
\item a \textbf{proof on paper};
\item the \textbf{physical meaning} (a blue box);
\item the \textbf{machine status} (a green box): unconditional, or conditional on named
  hypotheses, and the result of \texttt{\#print axioms} where it was run.
\end{enumerate}
Traps, dead ends and defects of the library are recorded in red boxes. Each chapter ends with a
\emph{Referee's checklist} (a table of its headline declarations) and with \emph{Notes and
references}. Section~\ref{ch01:sec:boxes} explains the boxes in more detail.

\subsection{How to use the book}

A reader who wants to verify the construction of $\Vnat$ can follow the \emph{Moonshine path} of
Section~\ref{ch01:sec:paths}; a reader interested in the library as a whole can follow the
\emph{general path}. In both cases the recommended way to use a chapter is the referee's way: read
the mathematical statement, then read the Lean statement \emph{before} reading the fidelity check,
and try to decode it yourself with the help of Sections~\ref{ch01:sec:lean}
and~\ref{ch01:sec:conventions}. Then read the fidelity check and see whether you agree. The book is
tied to a single commit: if the library has moved on, the script
\texttt{python3 scripts/textbook/sync.py check} lists every chapter whose quotations no longer match
the source (Section~\ref{ch01:sec:bookchecks}).

%======================================================================
\section{The Lean toolkit: reading Lean for mathematicians}\label{ch01:sec:lean}
%======================================================================

Lean~4 is a programming language and an interactive proof assistant based on dependent type
theory; Mathlib is its community library of formalised mathematics. The LeanVOA library is written
in Lean~4.30.0 against Mathlib v4.30.0. A reader of this book never needs to \emph{write} Lean; she
needs to \emph{read} statements. This section teaches the minimum, with short examples taken
verbatim from the library.

\subsection{Terms and types; propositions and data}

Every object in Lean is a \emph{term}, and every term has a \emph{type}, written \lean{x : X} (read:
``$x$ is a term of type $X$''). A type plays the role of the set to which an element belongs, with
one difference: the type of a term is part of the term and is always known to the system. Examples:
\lean{n : ℤ} is an integer, \lean{c : ℂ} a complex number, \lean{k : ℕ} a natural number (and
$\N$ \emph{includes} $0$). The expression \lean{(n : ℂ)} is the \emph{cast} of the integer $n$ into
$\C$; casts are sometimes printed as an upward arrow \lean{↑n}. \lean{𝒱 : Type u} declares $\choneV$
to be a type (``a set'') in universe level $u$; universes are a technical device to avoid Russell's
paradox and may be ignored. A function type is written with an arrow: \lean{f : ℤ → 𝒲} is a
function from $\Z$ to $\mathcal W$, applied by juxtaposition, \lean{f n}, so that \lean{f (-n - 1)}
means $f(-n-1)$. Functions of several variables are curried: \lean{g : ℤ → ℤ → 𝒲} is applied as
\lean{g m n}.

Statements are also types, of a special kind: they live in the universe \lean{Prop} of
propositions. A \emph{proof} of a statement \lean{P : Prop} is a term \lean{h : P}. So there are two
kinds of things in Lean, and a reader must always know which one she is looking at:
\begin{itemize}
\item \emph{data}, whose type lives in some \lean{Type u}: numbers, vectors, linear maps, vector
  spaces, vertex algebras;
\item \emph{proofs}, whose type is a proposition \lean{P : Prop}: the statement that a map is
  linear, that a family of subspaces is a direct-sum decomposition, that an identity holds.
\end{itemize}
Two proofs of the same proposition are regarded as equal (``proof irrelevance''), so a proof
carries no information beyond the truth of what it proves. Data, by contrast, can be chosen in
different ways, and the choice matters.

The logical connectives are written as in mathematics: \lean{P ∧ Q} (and), \lean{P ∨ Q} (or),
\lean{¬ P} (not), \lean{P → Q} (implies), \lean{P ↔ Q} (if and only if), \lean{∀ x, P x} (for
all), \lean{∃ x, P x} (there exists). \textbf{Quantifier order matters exactly as in mathematics.}
For instance \lean{∀ u v : 𝒱, ∃ N : ℤ, ∀ n : ℤ, N ≤ n → P u v n} says that a bound $N$ exists
for each pair $(u,v)$ and may depend on it; the statement with \lean{∃ N} written first would assert
a single uniform bound, which is much stronger. Implication associates to the right:
\lean{P → Q → R} means $P\Rightarrow(Q\Rightarrow R)$, i.e.\ ``if $P$ and $Q$, then $R$''. In a long
conjunction \lean{A ∧ B ∧ C} the connective also associates to the right; for reading purposes it is
simply the list of clauses $A$, $B$, $C$.


\textbf{Universes.} A declaration may say \lean{(𝒱 : Type u)}, \lean{(𝒲 : Type*)} or
\lean{Type _}. These are three ways of saying ``a type'' that differ only in how the universe level
is named: \lean{u} is a named level, \lean{Type*} introduces a fresh anonymous level, and \lean{_}
asks Lean to infer it. A referee can read all three as ``let $\choneV$ be a set''. The only place
where universes can matter for a statement is when two types are required to live in the
\emph{same} universe (a structure field such as \lean{carrier : Type u}); then a theorem quantifying
over all such structures quantifies over those in universe $u$, which for the mathematics of this
book loses nothing.

\subsection{Declarations: def, theorem, structure, class, instance}

A Lean file is a sequence of \emph{declarations}. The ones a reader meets are these.
\begin{itemize}
\item \lean{def name (arguments) : T := body} \emph{defines} a term of type \lean{T}. A definition
  proves nothing. If \lean{T} is \lean{Prop}, the definition introduces a new \emph{predicate}; much
  of the work of refereeing a library is checking that its predicates mean what their names
  suggest, because a theorem about a badly defined notion can be true and useless. This is why the
  book prints every definition that occurs in a headline theorem.
\item \lean{theorem name (arguments) : P := proof} asserts that the proposition \lean{P} holds for
  all values of the arguments. In this book we print the part up to and including the colon and the
  statement \lean{P}; the proof after \lean{:=} is usually omitted. (\lean{lemma} is a synonym.)
\item \lean{structure} introduces a new type whose terms are tuples of data and proofs ---
  Section~\ref{ch01:sec:structures} below.
\item \lean{class} is a structure whose terms Lean will \emph{search for automatically}; the
  standard examples are Mathlib's algebraic structures \lean{AddCommGroup 𝒱} and
  \lean{Module ℂ 𝒱}. The LeanVOA library declares essentially no classes of its own: a text search finds one
  declaration beginning with \lean{class}, \lean{VOA.Huang.TensorCat.FusionCategory}
  (Chapter~\ref{ch:L21}); its vertex-algebraic notions are ordinary structures, passed as explicit
  arguments.
\item \lean{instance} declares a term of a class, which the automatic search will then find.
\item \lean{abbrev} is a definition that Lean unfolds eagerly; for reading it is the same as
  \lean{def}. \lean{noncomputable} in front of a definition means that Lean cannot run it as a
  program (typically because it uses the axiom of choice); it has no logical cost for classical
  mathematics, but it means that numbers such as dimensions cannot simply be evaluated by the
  computer --- numerical evidence in the project is produced outside Lean and always labelled as
  evidence.
\item \lean{axiom name : P} \emph{postulates} \lean{P} without proof. The library contains none;
  until 19 October 2026 it contained one, \lean{VOA.pbw}, now a theorem
  (Section~\ref{ch01:sec:trust}).
\end{itemize}

\begin{example}[The mode convention]\label{ch01:ex:mode}
Here is the single most important definition of the library, in full.
\noindent File \leanpath{VOA/\allowbreak Foundations/\allowbreak FormalCalculus.lean}.
\begin{leancode}
def mode (f : FormalDist 𝒲) (m : ℤ) : 𝒲 := f (-m - 1)
\end{leancode}
Read it as follows. The name is \lean{mode} (in full, \lean{VOA.mode}, since the whole library sits
in the namespace \lean{VOA}). It has two explicit arguments, \lean{f} of type \lean{FormalDist 𝒲}
and \lean{m} of type \lean{ℤ}; the result has type \lean{𝒲}; and the value is \lean{f (-m - 1)}.
To understand it one must know what \lean{FormalDist 𝒲} is: it is the type \lean{ℤ → 𝒲} of all
functions from $\Z$ to $\mathcal W$, thought of as doubly infinite formal series
$f(z) = \sum_{k\in\Z} f(k)\,z^{k}$ (Chapter~\ref{ch:L01}, Definition~\ref{l01:def:formaldist});
the library also writes this type as \lean{𝒲⟦z^±⟧}. So \lean{f k} is the coefficient of $z^k$, and
\lean{mode f m} is the coefficient of $z^{-m-1}$, i.e.\ the $m$-th mode $f_{(m)}$ in
\begin{equation*}
  f(z) = \sum_{m\in\Z} f_{(m)}\, z^{-m-1}.
\end{equation*}
When \lean{f} is the vertex operator \lean{Y u} of a state $u$, \lean{mode (Y u) n} is the operator
$u_{(n)}$, whereas \lean{Y u k} (without \lean{mode}) is the coefficient of $z^k$, which is
$u_{(-k-1)}$. Both forms occur.

\emph{A worked example.} \lean{mode (Y u) 3} unfolds to \lean{Y u (-3 - 1)} $=$ \lean{Y u (-4)}, the
coefficient of $z^{-4}$ in $Y(u,z)$, and in $\sum_n u_{(n)}z^{-n-1}$ the power $z^{-4}$ carries
$u_{(3)}$, as it should. \emph{A check with the vacuum.} In a vertex algebra
$Y(u,z)\chonevac = u + z\,Tu + O(z^2)$, where $Tu = u_{(-2)}\chonevac$. The coefficient of $z^0$ is
\lean{Y u 0} $=$ \lean{mode (Y u) (-1)}, i.e.\ $u_{(-1)}$, and the library's creation axiom indeed
reads \lean{mode (Y u) (-1) vacuum = u}; the coefficient of $z^1$ is \lean{Y u 1} $=$
\lean{mode (Y u) (-2)}, i.e.\ $u_{(-2)}$, and the library defines the translation operator by
\lean{V.smul u (-2) V.vacuum}. The positive modes $u_{(n)}$, $n\ge0$, sit at negative powers
$z^{-n-1}$; they kill the vacuum, which is the statement that $Y(u,z)\chonevac$ has no pole at $z=0$.
The full discussion, with more calibrating examples, is Definition~\ref{l01:def:mode} of
Chapter~\ref{ch:L01}.
\end{example}

\begin{dictionary}
\textbf{How to read it.} Explicit arguments (in parentheses) are the inputs; the type after the last
colon is the output type. There is no hidden hypothesis: \lean{mode} is a plain re-indexing. The
variable \lean{𝒲} does not appear among the arguments because it was declared once for the whole
file by a \lean{variable} command (next subsection), which makes it an implicit argument.
\end{dictionary}

\begin{physics}
A chiral field of a two-dimensional conformal field theory is expanded in modes, and the
mathematician's convention $Y(u,z) = \sum u_{(n)}z^{-n-1}$ indexes the modes so that the exponent
does not depend on the conformal weight of $u$. For a field of weight $h$ the physicist writes
$\phi(z) = \sum\phi_m z^{-m-h}$, and then $u_{(n)} = \phi_{n-h+1}$; for the stress tensor ($h=2$)
this gives $L_m = \omega_{(m+1)}$. The definition itself is bookkeeping; the physics lies in the
axioms it is used to state (Chapter~\ref{ch:L03}).
\end{physics}

\begin{status}
A definition. \texttt{\#print axioms VOA.mode} reports that it depends on no axioms (run for this
chapter).
\end{status}

\begin{example}[An instance: the identity field]\label{ch01:ex:instOne}
\noindent File \leanpath{VOA/\allowbreak Core/\allowbreak VertexAlgebra.lean}.
\begin{leancode}
instance instOne [Zero 𝒲] [One 𝒲] : One (FormalDist 𝒲) := ⟨fun n => if n = 0 then 1 else 0⟩
\end{leancode}
This declares how the symbol \lean{1} is to be read for formal series: \lean{1} is the series whose
coefficient of $z^n$ is $1$ for $n=0$ and $0$ otherwise, i.e.\ the constant series $1$. The square
brackets \lean{[Zero 𝒲] [One 𝒲]} require that $\mathcal W$ have a zero and a one (for example
$\mathcal W = \End(V)$, where \lean{1} is the identity map). The angle brackets
\lean{⟨...⟩} build a term of the class \lean{One} from its single field. Why this matters: the vacuum
axiom of a vertex algebra is written \lean{Y vacuum = 1}, and with this instance it says
$Y(\chonevac,z) = \mathrm{id}_V$, i.e.\ $\chonevac_{(n)} = \delta_{n,-1}\,\mathrm{id}$. Had the
library used the \emph{pointwise} constant function $1$ on $\Z$ instead, \lean{1} would have meant
the formal delta function $\sum_n z^n$, and the vacuum axiom would have said something false for
every genuine vertex algebra. The library's choice is the correct one; the axiom is decoded in
Chapter~\ref{ch:L03}, Definition~\ref{l03:def:VA}.
\end{example}

\begin{dictionary}
An \lean{instance} has no name that the reader needs (here it is
\lean{VOA.FormalDist.instOne}); its effect is that every later occurrence of \lean{1} at the type
\lean{FormalDist 𝒲} means this series. A referee who sees \lean{1} or \lean{0} in a statement should
always ask \emph{which} instance gives it meaning.
\end{dictionary}

\begin{physics}
The identity field is the field of the vacuum state: inserting the vacuum at any point does
nothing. It is the unit of the operator product expansion.
\end{physics}

\begin{status}
A definition; \texttt{\#print axioms VOA.FormalDist.instOne} reports no axioms (run for this
chapter).
\end{status}

\subsection{Type synonyms, abbrev, and definitional unfolding}\label{ch01:sec:synonyms}

A \lean{def} whose value is a type creates a \emph{type synonym}: a new name for an existing type.
The library's most important one is the type of formal distributions.
\noindent File \leanpath{VOA/\allowbreak Foundations/\allowbreak FormalSeries.lean}.
\begin{leancode}
def FormalDist (𝒲 : Type*) : Type _ := ℤ → 𝒲
\end{leancode}
Mathematically \lean{FormalDist 𝒲} \emph{is} the set of functions $\Z\to\mathcal W$. But Lean
treats a \lean{def} as \emph{opaque} for the automatic search: the structures that Mathlib puts on
function types (pointwise addition, but also a pointwise \emph{multiplication}) are not inherited
automatically, and only the structures the library declares by hand are available. This is
deliberate: the pointwise product of two formal series in $z$ is \emph{not} their product as series
(and the series product does not exist in general), so the library keeps the misleading ring
structure away and gives \lean{FormalDist 𝒲} only its additive and $\C$-linear structure. A
side-effect, recorded in the source, is that general lemmas about functions do not apply
automatically, and the library has to restate trivialities such as the following.
\noindent File \leanpath{VOA/\allowbreak Core/\allowbreak Deformation.lean}.
\begin{leancode}
@[simp] theorem zero_dist_apply {𝒲 : Type*} [AddCommGroup 𝒲] (k : ℤ) :
    (0 : FormalDist 𝒲) k = 0 := rfl
\end{leancode}
Two pieces of syntax appear here. The proof \lean{rfl} (``reflexivity'') means that the two sides
are equal \emph{by unfolding definitions}: Lean checks the equation by computation. And
\lean{@[simp]} is an \emph{attribute} that registers the theorem as a rewriting rule for the
simplifier; it has no logical content. A theorem proved by \lean{rfl} is a \emph{definitional}
fact: it tells the referee exactly what a definition unfolds to. For instance, immediately after
\lean{mode}:
\noindent File \leanpath{VOA/\allowbreak Foundations/\allowbreak FormalCalculus.lean}.
\begin{leancode}
@[simp] theorem mode_apply (f : FormalDist 𝒲) (m : ℤ) : mode f m = f (-m - 1) := rfl
\end{leancode}
Such \lean{rfl}-lemmas are the safest possible statements: they cannot hide any mathematics. The
same holds for \lean{Iff.rfl}, the proof of \lean{P ↔ P}: a theorem \lean{A ↔ B := Iff.rfl} says that
the two statements are the same after unfolding definitions, so it is a membership or unfolding
lemma, not a theorem with content. Similarly \lean{funext h} proves \lean{f = g} from
\lean{h : ∀ n, f n = g n} (function extensionality). An example combining both:
\noindent File \leanpath{VOA/\allowbreak Foundations/\allowbreak FormalSeries.lean}.
\begin{leancode}
def LaurentPoly (𝒲 : Type*) [AddCommGroup 𝒲] [Module ℂ 𝒲] : Submodule ℂ (FormalDist 𝒲) where
  carrier := {f | (Function.support f).Finite}
\end{leancode}
\noindent File \leanpath{VOA/\allowbreak Foundations/\allowbreak FormalSeries.lean}.
\begin{leancode}
theorem mem_laurentPoly [AddCommGroup 𝒲] [Module ℂ 𝒲] {f : FormalDist 𝒲} :
    f ∈ LaurentPoly 𝒲 ↔ (Function.support f).Finite := Iff.rfl
\end{leancode}
A subspace is defined by giving its \lean{carrier} (the underlying set, here in set-builder notation
\lean{{f | ...}}) and proofs that it is closed under the operations (omitted here; the \lean{where}
syntax lists the fields of the structure being built). \lean{Function.support f} is the set
$\{n : f(n)\ne0\}$ and \lean{.Finite} the proposition that it is finite (Section~\ref{ch01:sec:sums});
\lean{mem_laurentPoly} merely unfolds the definition.

\textbf{Attributes.} The prefix \lean{@[...]} attaches an \emph{attribute} to a declaration: metadata
used by the automation, with no logical content. \lean{@[simp]} registers a rewrite rule for the
simplifier; \lean{@[ext]} registers an \emph{extensionality} lemma, i.e.\ a theorem ``two objects are
equal when their components agree'', which the tactic \lean{ext} then applies automatically;
\lean{@[inherit_doc]} copies a documentation string; \lean{@[reducible]} makes a definition unfold like
an \lean{abbrev}. Example:
\noindent File \leanpath{VOA/\allowbreak Foundations/\allowbreak FormalSeries.lean}.
\begin{leancode}
@[ext]
theorem ext {f g : FormalDist 𝒲} (h : ∀ n, f n = g n) : f = g := funext h
\end{leancode}
Two formal distributions are equal when all their coefficients are; since \lean{FormalDist} is opaque,
this has to be stated (and is proved by \lean{funext}). Likewise \lean{Aut.ext}
(\leanpath{VOA/\allowbreak Core/\allowbreak Automorphism.lean}) says that two automorphisms are equal when their
linear maps are: the proof fields carry no information.

\textbf{\lean{inferInstanceAs}.} Since an opaque type synonym inherits no structure, the library
transfers the structures it \emph{does} want by hand:
\noindent File \leanpath{VOA/\allowbreak Foundations/\allowbreak FormalSeries.lean}.
\begin{leancode}
instance instAddCommGroup [AddCommGroup 𝒲] : AddCommGroup (FormalDist 𝒲) :=
  inferInstanceAs (AddCommGroup (ℤ → 𝒲))
\end{leancode}
\lean{inferInstanceAs (C X)} means ``the instance of \lean{C X} that the automatic search finds'', used
here at the synonym. So the addition on formal distributions is exactly the pointwise addition of
functions $\Z\to\mathcal W$; the same is done for scalar multiplication, and nothing is done for
multiplication. A referee reads \lean{inferInstanceAs} as ``the evident structure, transported''.

\textbf{\lean{abbrev} versus \lean{def}.} An \lean{abbrev} is a definition that Lean unfolds
\emph{eagerly}, in particular during the automatic instance search, so that an \lean{abbrev} of a
type inherits all the structure of the type it abbreviates. Logically the two are the same: the
abbreviated type and its unfolding are definitionally equal, and no statement changes meaning. The
difference is practical, and it can be large. An example:
\noindent File \leanpath{VOA/\allowbreak Moonshine/\allowbreak ExtraspecialModule.lean}.
\begin{leancode}
noncomputable abbrev ExtraT : Type := Rep (cocMat L) (isUnit_det_cocMat L hEven hU)
\end{leancode}
Because this abbreviation is unfolded during instance search, elaborating even the scalar action
on its elements times out, and downstream files use the unfolded type \lean{Rep (cocMat L) _}
directly (Chapter~\ref{ch:06}). For a referee the lesson is only that the \emph{same} mathematical
object can appear under several names in different files; the fidelity checks say when two names are
definitionally equal. Conversely the carrier of the Moonshine module is an \lean{abbrev},
\lean{abbrev Vnat : Type := SkewVP × TP} (Section~\ref{ch01:sec:linalg}), so that it inherits the
vector-space structure of the product.

\begin{warning}
\textbf{How type synonyms can mislead.} If a type synonym \emph{did} inherit a structure the
authors did not intend (for \lean{FormalDist}, the pointwise product), a statement using \lean{*} on
it would silently mean the wrong product. The library avoids this for \lean{FormalDist} by making it
opaque. A referee who sees an algebraic operation on a type synonym should ask which instance gives
it meaning --- the same question as for \lean{1} in Example~\ref{ch01:ex:instOne}.
\end{warning}

\subsection{Arguments: explicit, implicit, instance}

Lean has three kinds of brackets for arguments, and all three are \emph{universally quantified} for
reading purposes.
\begin{itemize}
\item \lean{(x : X)}: an \emph{explicit} argument, written out when the declaration is used.
\item \lean{{x : X}}: an \emph{implicit} argument, which Lean infers from the other arguments (for
  example the type $\choneV$ of a vector \lean{v : 𝒱}). In a statement, \lean{∀ {n : ℤ} (v : 𝒱), ...}
  means ``for all $n$ and $v$''; the braces only say that $n$ need not be written when the statement
  is used.
\item \lean{[C X]}: an \emph{instance} argument, a structure of class \lean{C} on \lean{X} which Lean
  finds by the automatic search.
\item \lean{⦃x : X⦄}: a \emph{strict implicit} argument. Like \lean{{x : X}} it is inferred, but only
  once a later explicit argument is supplied. For reading it is again ``for all $x$''. Example:
\end{itemize}
\noindent File \leanpath{VOA/\allowbreak Modularity/\allowbreak SuperTorusRecursion.lean}.
\begin{leancode}
def IsParityHomogeneousMap {X Y : Type*} [AddCommGroup X] [Module ℂ X]
    [AddCommGroup Y] [Module ℂ Y] (P : ParityGrading X) (Q : ParityGrading Y)
    (p : ZMod 2) (f : X →ₗ[ℂ] Y) : Prop :=
  ∀ (q : ZMod 2) ⦃x : X⦄, x ∈ P.parity q → f x ∈ Q.parity (p + q)
\end{leancode}
which reads: $f$ is homogeneous of parity $p$ if for every $q$ and every $x$ of parity $q$, $f(x)$ has
parity $p+q$ (addition in \lean{ZMod 2}).

\textbf{Other declaration syntax.} \lean{inductive} defines a type by listing its constructors (for
example a type with finitely many named elements, or a type of formal words); \lean{deriving
DecidableEq} asks Lean to generate the evident instance. \lean{open X in} makes the short names of
namespace \lean{X} available for one declaration, and \lean{scoped} notation is notation available
only after \lean{open}. A name prefixed by \lean{_root_.}, as in \lean{_root_.Even}, refers to the
global declaration rather than to one of the same name inside the current namespace. Tactic names
(\lean{simp}, \lean{simpa}, \lean{rw}, \lean{exact}, \lean{omega}, \lean{norm_num}, \lean{ring}, \lean{ext})
and names of Mathlib lemmas such as \lean{smul_add} occur only inside proofs; a reader of statements
can ignore them.
Almost every file of the library begins with a line such as
\noindent File \leanpath{VOA/\allowbreak Core/\allowbreak VertexAlgebra.lean}.
\begin{leancode}
variable {𝒱 : Type u} [AddCommGroup 𝒱] [Module ℂ 𝒱]
\end{leancode}
which reads: \emph{let $\choneV$ be a complex vector space}. \lean{AddCommGroup 𝒱} is an abelian
group structure and \lean{Module ℂ 𝒱} a compatible $\C$-module structure. The \lean{variable}
command adds these arguments to every subsequent declaration of the file that mentions $\choneV$, so
they do not appear in the printed statement but are nevertheless hypotheses of it. They are harmless
here --- they say only that $\choneV$ is a complex vector space --- but a referee must remember that
a statement can have arguments that are declared further up the file. The fidelity checks of the
book always list them.


A \lean{variable} line may also introduce an \emph{explicit} argument, for instance
\lean{(W :}\allowbreak\ \lean{VertexOperatorAlgebra 𝒱)}. Then every definition of the file that uses \lean{W} takes
it as its first explicit argument although it is not printed in the definition itself:
\noindent File \leanpath{VOA/\allowbreak Moonshine/\allowbreak ZhuGenFiltration.lean}.
\begin{leancode}
variable {𝒱 : Type u} [AddCommGroup 𝒱] [Module ℂ 𝒱] (W : VertexOperatorAlgebra 𝒱)
noncomputable def Fl (n : ℤ) : Submodule ℂ 𝒱 := ⨆ (k : ℤ) (_ : k < n), W.grading k
\end{leancode}
Here \lean{Fl} is really a function of two explicit arguments, \lean{Fl W n}, the span
$\bigoplus_{k<n}W_k$ of all graded pieces of weight below $n$ (the notation \lean{⨆} is explained in
Section~\ref{ch01:sec:linalg}).

\textbf{Typeclass search, and why its depth matters.} When a declaration needs an instance argument
\lean{[C X]}, Lean searches the registered instances, possibly chaining several (``$\choneV$ is a
$\C$-module, hence an additive group, hence \dots''). The search is bounded by options, one of which,
\lean{maxSynthPendingDepth}, the library sets globally in its build configuration:
\begin{plaincode}
[leanOptions]
maxSynthPendingDepth = 3
\end{plaincode}
(an excerpt of \texttt{lakefile.toml}, shown as plain text because it is not Lean source). Some
files need this larger depth to elaborate. If a file is compiled \emph{without} it --- which is what
happens with a bare \texttt{lake env lean}, since that command does not read the build options ---
the search can fail, and the elaborator may then replace the failing piece by \lean{sorryAx} and
report an error. This is the origin of the ``audit trap'' of Section~\ref{ch01:sec:trust}: axiom
audits must be run through \texttt{lake build} or with the option passed explicitly. A failed
instance search never produces a \emph{wrong} accepted theorem; it produces an error (or a
\lean{sorryAx} that \texttt{\#print axioms} reveals).

\subsection{Structures: a bundle of data and proofs}\label{ch01:sec:structures}

A \lean{structure} defines a mathematical object as a tuple of data together with axioms --- exactly
like ``a group is a set with a multiplication, a unit and inverses, such that \dots''. Each component
is a \emph{field} \lean{name : type}. Fields whose type is data are the data; fields whose type is a
proposition are the axioms, and a term of the structure must supply a proof of each of them.

\begin{example}[Automorphisms of a VOA]\label{ch01:ex:aut}
\noindent File \leanpath{VOA/\allowbreak Core/\allowbreak Automorphism.lean}.
\begin{leancode}
structure Aut (W : VertexOperatorAlgebra 𝒱) where
  toLinearEquiv : 𝒱 ≃ₗ[ℂ] 𝒱
  map_smul' : ∀ (u : 𝒱) (n : ℤ) (v : 𝒱),
    toLinearEquiv (W.smul u n v) = W.smul (toLinearEquiv u) n (toLinearEquiv v)
  map_vacuum : toLinearEquiv W.vacuum = W.vacuum
  map_omega : toLinearEquiv W.conformalVector = W.conformalVector
\end{leancode}
(documentation comments removed). A term \lean{α : Aut W} is a quadruple: one piece of data, the
linear isomorphism \lean{toLinearEquiv} $:\choneV\to\choneV$, and three proofs: that it preserves
every $n$-th product, $\alpha(u_{(n)}v) = (\alpha u)_{(n)}(\alpha v)$; that it fixes the vacuum;
and that it fixes the conformal vector. The components are accessed by \emph{projection}, written
with a dot: \lean{α.toLinearEquiv}, \lean{α.map_vacuum}. The same dot notation is used for functions
defined in a namespace: if \lean{W : VertexOperatorAlgebra 𝒱} then \lean{W.smul u n v} is the
function \lean{PreVertexAlgebra.smul} applied to (the underlying pre-vertex algebra of) \lean{W} and to
\lean{u n v}; it is the $n$-th product $u_{(n)}v$ --- the name \lean{smul} refers to ``product'', not
to scalar multiplication. The mathematics of automorphisms (that they form a group, preserve the
grading, and so on) is in Chapter~\ref{ch:L04}, Definition~\ref{l04:def:aut}.
\end{example}

\begin{dictionary}
\textbf{Decoding the symbols.} \lean{𝒱 ≃ₗ[ℂ] 𝒱} is the type of $\C$-linear isomorphisms
$\choneV\to\choneV$, bundled together with their inverses. The quantifier in \lean{map_smul'} is over
all $u$, $n$, $v$, with no restriction. The structure takes the VOA \lean{W} as an explicit
parameter, so \lean{Aut W} is a type that depends on \lean{W}. Nothing requires the grading to be
preserved: that is a \emph{theorem} (Chapter~\ref{ch:L04}), proved from $\alpha\omega = \omega$.
\end{dictionary}

\begin{physics}
An automorphism of the chiral algebra is a global internal symmetry of the conformal field theory
that commutes with conformal transformations (it fixes the stress tensor). For $\Vnat$ the
automorphism group is the Monster (a theorem of Frenkel, Lepowsky and Meurman, not formalized in
the library).
\end{physics}

\begin{status}
A definition; \texttt{\#print axioms VOA.Aut} lists \lean{[propext, Classical.choice, Quot.sound]}
(run for this chapter).
\end{status}

\textbf{The most dangerous failure of a structure.} A structure whose axioms are contradictory has no
terms at all, and then every theorem of the form ``for every $W$ of this structure, \dots'' is
vacuously true. The book checks for this explicitly; the library has had several such cases
(Section~\ref{ch01:sec:failures}, item (F1)). One, a twisted-module axiom quantified over too many
indices, is recorded and repaired in Chapter~\ref{ch:06}; the others are recorded, unrepaired, in
Chapters~\ref{ch:L12} and~\ref{ch:L17}. The opposite failure is a structure satisfied by a trivial object: for
instance, the zero space is a vertex operator algebra of every central charge, which is why
headline results assert simplicity.

A structure may \lean{extend} another: \lean{structure B extends A} means that a term of \lean{B} has
all the fields of \lean{A} plus new ones, and \lean{b.toA} is its underlying \lean{A}. This is how a
vertex operator algebra is built, in layers, on top of the library's basic structure \lean{PreVertexAlgebra}
(Example~\ref{ch01:ex:VOA}).

\subsection{The Mathlib vocabulary at a glance}

The following Mathlib notions occur constantly. They are standard mathematics; we fix the notation
and record the conventions that a referee must know. This is an overview; each item is treated with
library examples in Sections~\ref{ch01:sec:choice}--\ref{ch01:sec:algebra}, and the glossary of
Section~\ref{ch01:sec:glossary} lists everything in one table.
\begin{itemize}
\item \lean{𝒱 →ₗ[ℂ] 𝒲}: the $\C$-linear maps $\choneV\to\mathcal W$, applied like functions,
  \lean{φ v}. \lean{𝒱 ≃ₗ[ℂ] 𝒲}: the $\C$-linear isomorphisms.
\item \lean{Module.End ℂ 𝒱}, written \lean{End ℂ 𝒱} in the library: the algebra $\End_\C(\choneV)$.
  In it \lean{a * b} is the composition $a\circ b$, \lean{1} is the identity, and \lean{⁅a, b⁆} is
  the commutator $ab-ba$.
\item \lean{Submodule ℂ 𝒱}: the linear subspaces of $\choneV$, ordered by inclusion, with least
  element \lean{⊥} $=\{0\}$ and greatest element \lean{⊤} $=\choneV$. Membership is \lean{v ∈ S};
  \lean{ℂ ∙ v} is the line $\C v$; \lean{Submodule.span ℂ s} is the span of a set $s$.
\item \lean{↥S}: a subspace \lean{S : Submodule ℂ 𝒱} is a set of vectors; \lean{↥S} is the same
  subspace regarded as a vector space in its own right (its elements are pairs of a vector and a
  proof that it lies in $S$). The arrow is a \emph{coercion}: Lean inserts it automatically when a
  subspace is used where a type is expected, and it is usually invisible in the source but visible
  in Lean's printed output.
\item \lean{Module.finrank ℂ X}: the dimension of $X$ as a natural number. \textbf{Caution:} by
  Mathlib's convention the \lean{finrank} of an \emph{infinite}-dimensional space is $0$. So
  \lean{finrank ℂ X = 196884} implies that $X$ is finite-dimensional, but \lean{finrank ℂ X = 0}
  does not imply $X = 0$. This convention has bitten the library once: its ``fusion coefficient'' is
  identically zero for exactly this reason (Chapter~\ref{ch:L12}, Theorem~\ref{l12:thm:fusionzero}).
  \lean{FiniteDimensional ℂ X} is the proposition that $X$ is finite-dimensional.
\item \lean{DirectSum.IsInternal S}, for a family \lean{S : ℤ → Submodule ℂ 𝒱}: the proposition
  that the natural map $\bigoplus_n S_n\to\choneV$ is an isomorphism, i.e.\
  $\choneV = \bigoplus_n S_n$ internally. This is how the library says that a space is graded.
\item \lean{LinearMap.BilinForm ℂ 𝒱}: bilinear forms $\choneV\times\choneV\to\C$, applied as
  \lean{B a b}.
\item \lean{PowerSeries ℂ}: formal power series $\sum_{n\ge0}a_nq^n$; \lean{PowerSeries.coeff n f}
  is the coefficient $a_n$ (printed \lean{(PowerSeries.coeff n) f}).
\item \lean{∑ i ∈ Finset.range N, f i} is the finite sum $\sum_{i=0}^{N-1}f(i)$. \lean{∑ᶠ i, f i}
  (``\lean{finsum}'') is the sum of a family with finite support, and \textbf{by convention it is $0$
  when the support is infinite}. Wherever the library uses \lean{∑ᶠ}, finiteness of the support must
  be checked separately; the chapters point this out where it matters (see
  Chapter~\ref{ch:L01}, warning~\ref{l01:warn:finsum}).
\item \lean{ZMod n}: the integers modulo $n$; \lean{ZMod 2} $=\{0,1\}$ is the parity group, used for
  the even/odd grading of superalgebras and for the $\Z/2$ orbifold.
\end{itemize}

\begin{example}[Reading a structure header: vertex operator algebras]\label{ch01:ex:VOA}
The central structure of the library is the following.
\noindent File \leanpath{VOA/\allowbreak Core/\allowbreak VertexOperatorAlgebra.lean}.
\begin{leancode}
structure GradedVertexOperatorAlgebra (𝒱 : Type u) [AddCommGroup 𝒱] [Module ℂ 𝒱]
    extends PreVertexAlgebra 𝒱 where
  conformalVector : 𝒱
  conformal_mem : conformalVector ∈ grading 2
  centralCharge : ℂ
  virasoro : ∀ m n : ℤ,
      ⁅mode (Y conformalVector) (m + 1), mode (Y conformalVector) (n + 1)⁆
        = ((m : ℂ) - n) • mode (Y conformalVector) (m + n + 1)
          + (if m + n = 0 then (((m : ℂ) ^ 3 - m) / 12 * centralCharge) • (1 : End ℂ 𝒱)
             else 0)
  L0_grading : ∀ {n : ℤ} (v : 𝒱), v ∈ grading n → mode (Y conformalVector) 1 v = (n : ℂ) • v
  translation : ∀ v : 𝒱, Y (mode (Y conformalVector) 0 v) = deriv (Y v)
  grading_bounded_below : ∃ N : ℤ, ∀ n : ℤ, n < N → grading n = ⊥

structure VertexOperatorAlgebra (𝒱 : Type u) [AddCommGroup 𝒱] [Module ℂ 𝒱]
    extends GradedVertexOperatorAlgebra 𝒱 where
  vacuum_grade : grading 0 = ℂ ∙ vacuum
  positivity : ∀ n : ℤ, n < 0 → grading n = ⊥
  grading_bounded_below := ⟨0, positivity⟩

def L (W : VertexOperatorAlgebra 𝒱) (n : ℤ) : End ℂ 𝒱 :=
  mode (W.Y W.conformalVector) (n + 1)
\end{leancode}
(documentation comments removed). We read it only at the level of ``what kind of thing is each
line''; the mathematics is in Chapter~\ref{ch:L03}, Definition~\ref{l03:def:VOA}.
\begin{itemize}
\item The headers: \lean{GradedVertexOperatorAlgebra 𝒱} is a type depending on a complex vector space
  $\choneV$; \lean{extends PreVertexAlgebra 𝒱} means that it has all the fields of
  \lean{PreVertexAlgebra} --- a grading \lean{grading : ℤ → Submodule ℂ 𝒱}, its directness, the
  vertex operator \lean{Y : 𝒱 →ₗ[ℂ] (End ℂ 𝒱)⟦z^±⟧}, the \lean{vacuum} and Gannon's axioms VA1--VA5
  --- which is why \lean{grading}, \lean{Y} and \lean{vacuum} can be used in the new fields
  without being declared here. In the same way \lean{VertexOperatorAlgebra 𝒱} \lean{extends
  GradedVertexOperatorAlgebra 𝒱}, so a VOA has every field of both structures.
\item The layers. \lean{PreVertexAlgebra} (Chapter~\ref{ch:L03}) is \emph{not} a vertex algebra in the
  standard sense: it has no translation-covariance axiom, and the Borcherds identity can fail for it.
  The library's \lean{VertexAlgebra} is a \lean{PreVertexAlgebra} plus that axiom.
  \lean{GradedVertexOperatorAlgebra} is the vertex operator algebra of Frenkel--Lepowsky--Meurman and
  Lepowsky--Li: finite-dimensional graded pieces, the grading truncated from below
  (\lean{grading_bounded_below}), Virasoro, $L_0$-grading and $L_{-1}$-derivative.
  \lean{VertexOperatorAlgebra} adds the two \emph{CFT-type} fields \lean{vacuum_grade} and
  \lean{positivity}; the line \lean{grading_bounded_below := ⟨0, positivity⟩} fills the parent's truncation
  field automatically, so a construction of a VOA need not supply it. Every VOA is a vertex algebra
  (\lean{VertexOperatorAlgebra.toVertexAlgebra}), because \lean{translation} gives covariance.
\item Data: \lean{conformalVector} (a vector $\omega$) and \lean{centralCharge} (a number
  $c\in\C$).
\item Axioms: every other field is a proposition. \lean{conformal_mem} says $\omega\in V_2$.
  \lean{virasoro} is an equation in \lean{End ℂ 𝒱} quantified over all $m,n\in\Z$, with the
  commutator bracket \lean{⁅·,·⁆}, a cast \lean{(m : ℂ)}, scalar multiplication \lean{•}, and an
  \lean{if ... then ... else} for the central term: by the mode convention of
  Example~\ref{ch01:ex:mode}, \lean{mode (Y conformalVector) (m + 1)} is $\omega_{(m+1)} = L_m$, so
  this is the Virasoro relation. \lean{L0_grading} has an implicit weight \lean{{n : ℤ}};
  \lean{translation} uses the formal derivative \lean{deriv}; \lean{vacuum_grade} uses the line
  \lean{ℂ ∙ vacuum}; \lean{positivity} and \lean{grading_bounded_below} use \lean{⊥}.
\item After the structure, \lean{def L} introduces the abbreviation \lean{W.L n} for
  $L_n = \omega_{(n+1)}$.
\end{itemize}
The library also has a separate structure \lean{VOA.NearVOA} in the same file, with the same
Virasoro, grading and translation axioms but without finite-dimensionality of the graded pieces and
without the conditions $V_0 = \C\chonevac$, $V_{n<0}=0$; every VOA gives one by
\lean{VertexOperatorAlgebra.toNearVOA}.
\end{example}

\begin{dictionary}
\textbf{What to check when reading a structure.} (a) Which fields are data and which are axioms.
(b) The quantifiers of each axiom, including implicit ones in braces. (c) Every convention in an
axiom --- here the shift $L_m = \omega_{(m+1)}$, which would make the Virasoro axiom false for every
genuine VOA if it were off by one. (d) Whether the axioms can be satisfied by something non-trivial,
and whether they are satisfied by something trivial (the zero space is a VOA of every central
charge). (e) What the parent structure contributes. The full fidelity check of this structure ---
including the important point that \lean{VertexOperatorAlgebra} is of \emph{CFT type} by definition, which is
stronger than the ``grading bounded below'' of Frenkel--Lepowsky--Meurman (that general notion is
\lean{GradedVertexOperatorAlgebra}) --- is in Chapter~\ref{ch:L03}.
\end{dictionary}

\begin{warning}
\textbf{Correction: a bare \lean{PreVertexAlgebra} has no translation covariance.} An earlier version
of this chapter called the library's \lean{PreVertexAlgebra} ``stronger'' than the textbook
definitions of Kac and Frenkel--Ben-Zvi. That was wrong. \lean{PreVertexAlgebra} has a grading with
finite-dimensional pieces but \emph{no} translation-covariance axiom, and without it the Borcherds
(Jacobi) identity does not follow: Chapter~\ref{ch:L03} gives an explicit three-dimensional
\lean{PreVertexAlgebra} on which both the Borcherds identity and $Y(Tu,z) = \partial_zY(u,z)$ fail,
machine-checked and axiom-clean (a copy is in the library as
\lean{VOA.TranslationCovariantCounterexample.VA}).
Every Borcherds, skew-symmetry and duality theorem of the library is therefore stated for
\lean{VertexOperatorAlgebra}, whose axiom \lean{translation} restores covariance, and every
statement about $\Vnat$ is about \lean{VertexOperatorAlgebra}; the headline theorem is unaffected.
The defect is catalogued in Chapter~\ref{ch:15}, Section~\ref{ch15:sec:debt}; it led the library to rename the structure (\lean{PreVertexAlgebra}) and to reserve \lean{VertexAlgebra} for the translation-covariant notion.
\end{warning}

\begin{physics}
A vertex operator algebra is the chiral algebra of a two-dimensional conformal field theory together
with its stress tensor $T(z) = Y(\omega,z) = \sum L_nz^{-n-2}$; the central charge $c$ is the
conformal anomaly, the grading by $L_0$ is the grading by conformal weight, and the axiom
\lean{translation} is the Ward identity saying that $L_{-1}$ generates translations.
Section~\ref{ch01:sec:physics} gives the picture in a little more detail.
\end{physics}

\begin{status}
A definition. \texttt{\#print axioms VOA.VertexOperatorAlgebra} lists
\lean{[propext, Classical.choice, Quot.sound]} (run for this chapter).
\end{status}

\subsection{Existence, emptiness and choice}\label{ch01:sec:choice}

\textbf{\lean{∃} and \lean{Nonempty}.} \lean{∃ x : X, P x} is a proposition; a proof of it does not
come with a usable witness unless one invokes choice. For a type \lean{X},
\lean{Nonempty X} is the proposition ``$X$ has at least one element'', and \lean{IsEmpty X} the
proposition ``$X$ has no element''. Theorems of the form \lean{¬ Nonempty X} are how the library
records that a notion is empty --- which is sometimes exactly the warning a referee needs:
\noindent File \leanpath{VOA/\allowbreak Modules/\allowbreak ZhuVermaAnArchZhu534.lean}.
\begin{leancode}
theorem not_nonempty_inducedAssignment (W : VertexOperatorAlgebra 𝒱) :
    ¬ Nonempty (InducedAssignment W) := by
\end{leancode}
Every theorem that takes a term of \lean{InducedAssignment W} as a hypothesis is therefore vacuous
(Chapter~\ref{ch:L12}): the structure is uninhabited, so such a theorem proves nothing about any module.

\textbf{\lean{Function.Injective}.} \lean{Function.Injective f} is the proposition
$\forall a\,b,\ f\,a = f\,b\to a = b$. (Surjective and bijective are \lean{Function.Surjective},
\lean{Function.Bijective}.) It appears, for instance, in the PBW principle \lean{VOA.pbw} (Chapter~\ref{ch:L02}).

\textbf{Choice: \lean{.choose} and \lean{Classical.choose}.} From a proof \lean{h : ∃ x, P x}
Lean can extract a witness \lean{h.choose} (also written \lean{Classical.choose h}) together with
the proof \lean{h.choose_spec : P h.choose}. The witness is \emph{some} element satisfying $P$,
selected by the axiom of choice; nothing is known about it beyond $P$. Two library examples:
\noindent File \leanpath{VOA/\allowbreak Core/\allowbreak Automorphism.lean}.
\begin{leancode}
def IsLocNilp (D : End ℂ 𝒱) : Prop := ∀ v : 𝒱, ∃ N : ℕ, (D ^ N) v = 0

noncomputable def locNilpIdx {D : End ℂ 𝒱} (hD : IsLocNilp D) (v : 𝒱) : ℕ := (hD v).choose
\end{leancode}
\noindent File \leanpath{VOA/\allowbreak Moonshine/\allowbreak SigmaAssemble.lean}.
\begin{leancode}
noncomputable def s0 : ZMod 2 := (seed_in_frame seedExists).choose
\end{leancode}
The first chooses, for a locally nilpotent operator $D$ and a vector $v$, \emph{some} $N$ with
$D^Nv = 0$ (not necessarily the least one). The second chooses, in the Moonshine construction, one of
the two sectors in which a ``seed'' vector exists. Such definitions are \lean{noncomputable}. Mathlib
provides choice-based constructions of the same kind, the commonest in this library being a basis of
a free module:
\noindent File \leanpath{VOA/\allowbreak Foundations/\allowbreak LatticeCocycle.lean}.
\begin{leancode}
noncomputable def latticeBasis :
    Basis (Fin (Fintype.card (Module.Free.ChooseBasisIndex ℤ L.carrier))) ℤ L.carrier :=
  (Module.Free.chooseBasis ℤ L.carrier).reindex (Fintype.equivFin _)
\end{leancode}
\lean{Module.Free.chooseBasis ℤ L.carrier} is \emph{an arbitrary} $\Z$-basis of the lattice $L$,
re-indexed by \lean{Fin n} $=\{0,\dots,n-1\}$ (\lean{Fintype.card} is the number of elements of a
finite type).

\begin{warning}
\textbf{How choice can make a statement weaker than it looks.} A definition built from a chosen
witness depends on the choice, and a theorem about it is a theorem about \emph{that} choice. This is
harmless when the theorem holds for every choice --- the FLM cocycle built from
\lean{latticeBasis} satisfies the cocycle identities for every basis (Chapters~\ref{ch:L02}
and~\ref{ch:03}) --- and it is what the referee must check. A statement such as ``\lean{locNilpIdx}
is minimal'' would be unprovable, since nothing about the chosen $N$ beyond $D^Nv = 0$ is known. In
the Moonshine construction the choices (\lean{s0}, the seed vector, the basis) are all of this
harmless kind: the headline theorem is an existence statement, and its conclusion does not refer to
them (Chapter~\ref{ch:14}).
\end{warning}

\subsection{Total functions and junk values}\label{ch01:sec:junk}

Every Lean function is \emph{total}: it returns a value on every input, including inputs where the
mathematical operation is undefined. The value returned there is a \emph{junk value}, fixed by
convention. The ones that matter:
\begin{itemize}
\item \textbf{division by zero}: \lean{x / 0 = 0} in $\C$, $\Q$, $\mathbb R$; in $\N$ and $\Z$
  division is division with remainder rounded down, so \lean{7 / 2 = 3};
\item \textbf{natural subtraction is truncated}: for \lean{m n : ℕ}, \lean{m - n = 0} when $m<n$;
\item \lean{Int.toNat}: \lean{k.toNat} is $k$ for $k\ge0$ and $0$ for $k<0$;
\item \lean{Module.finrank} of an infinite-dimensional space is $0$, and \lean{∑ᶠ} of an infinitely
  supported family is $0$ (Sections~\ref{ch01:sec:linalg} and~\ref{ch01:sec:sums});
\item library functions that are only meaningful on some inputs return $0$ elsewhere.
\end{itemize}
Junk values are harmless when the statement only uses the function where it is meaningful, and
the referee's job is to check that. Three library examples.
\noindent File \leanpath{VOA/\allowbreak Foundations/\allowbreak Substitution.lean}.
\begin{leancode}
noncomputable def cbz (n k : ℤ) : ℂ := if 0 ≤ k then gBinom n k.toNat else 0
\end{leancode}
The binomial $\binom nk$ with an \emph{integer} lower index: \lean{k.toNat} is used only under the
guard \lean{0 ≤ k}, and for $k<0$ the value is \emph{defined} to be $0$. This junk value is
mathematically the right one: it lets the library write the Borcherds identity as a sum over all
$i\in\Z$ (Chapter~\ref{ch:L01}).
\noindent File \leanpath{VOA/\allowbreak Moonshine/\allowbreak TwistedHeisenberg.lean}.
\begin{leancode}
noncomputable def modeT (r : ℕ) (i : Fin r) (k : ℤ) : Module.End ℂ (FockT r) :=
  if k % 2 = 0 then 0
  else if k < 0 then mulXT (varT i k)
  else ((k : ℂ) / 2) • pderivT (varT i k)
\end{leancode}
This is the \emph{odd-integer encoding} of half-integer modes: the twisted boson has modes
$h_i(r)$ only for $r\in\frac12+\Z$, and the library indexes $h_i(k/2)$ by the odd integer $k$. For
even $k$ (\texttt{k \% 2 = 0}, where \texttt{\%} is the remainder) there is no mode, and the junk value
\lean{0} is returned. The factor \lean{(k : ℂ) / 2} is honest complex division. A statement about
\lean{modeT} for \emph{all} \lean{k} therefore also says something (usually trivial) about the junk
values; the fidelity checks of Chapter~\ref{ch:06} note where hypotheses ``$k$ odd'' are needed.
\noindent File \leanpath{VOA/\allowbreak Moonshine/\allowbreak ExtraspecialModule.lean}.
\begin{leancode}
theorem finrank_extraT : Module.finrank ℂ (ExtraT L hEven hU) = 2 ^ (L.rank / 2) := by
\end{leancode}
Here \lean{L.rank / 2} is division of natural numbers, rounded down. The formula $\dim T_L =
2^{\mathrm{rank}\,L/2}$ is meant for even rank, and the library proves separately that an even
unimodular lattice has even rank (\lean{even_rank_of_unimodular}, same file), so no rounding occurs
in the cases used.

\begin{warning}
\textbf{How junk values make statements vacuous or wrong.} A conclusion
\lean{Module.finrank ℂ X = 0} is satisfied by every infinite-dimensional $X$; a conclusion
\lean{∑ᶠ i, f i = 0} by every infinitely supported $f$; an equation between natural numbers involving
\lean{m - n} may hold only because both sides are truncated. The library has one real casualty: its
``fusion coefficient'' \lean{VOA.fusionCoeff} is a \lean{finrank} of a space that is always either
$0$ or infinite-dimensional, so it is identically $0$ (Chapter~\ref{ch:L12},
Theorem~\ref{l12:thm:fusionzero}). Every statement
whose content is a value of \lean{VOA.fusionCoeff} is therefore vacuous or false as a fusion rule.
\end{warning}

\subsection{Numbers, casts, ZMod, Fin, and decidable conditions}\label{ch01:sec:numbers}

\textbf{Casts.} $\N$, $\Z$, $\Q$, $\mathbb R$, $\C$ are different types, and a number is moved between
them by a cast: \lean{(n : ℤ)} for $n\in\N$, \lean{(m : ℂ)} for $m\in\Z$, printed \lean{↑n} by Lean.
Casts are injective ring homomorphisms, so an equation of casts means the evident equation, with two
cautions: a subtraction or division \emph{inside} a cast is computed in the source type (so
\lean{((m - n : ℕ) : ℤ)} is truncated, whereas \lean{(m : ℤ) - n} is not), and the position of the
cast matters for that reason. The headline theorem casts the natural number \lean{finrank} to $\C$ to
compare it with a power-series coefficient (Section~\ref{ch01:sec:headline}).

\textbf{\lean{ZMod n} and \lean{Fin n}.} \lean{ZMod n} is the ring $\Z/n\Z$; \lean{ZMod 2} is the
parity group, used for super gradings, for the two sectors of an orbifold, and in coding theory.
\lean{Fin n} is the finite type $\{0,1,\dots,n-1\}$, used as an index set: \lean{i : Fin r} in
\lean{modeT} above is a colour index $1\le i\le r$ (numbered from $0$). An element of \lean{Fin n}
is a natural number together with a proof that it is less than $n$.

\textbf{Decidable conditions.} An expression \lean{if c then a else b} requires Lean to be able to
decide \lean{c}; for equalities of integers, inequalities and parity this is automatic (the relevant
instances are called \lean{Decidable} and \lean{DecidableEq}), and with classical logic every
proposition is decidable. For reading purposes \lean{if c then a else b} is the case distinction it
looks like. Example:
\noindent File \leanpath{VOA/\allowbreak Foundations/\allowbreak Expansion.lean}.
\begin{leancode}
def negOnePow (n : ℤ) : ℂ := if Even n then 1 else -1
\end{leancode}
which is $(-1)^n$ (\lean{Even n} is the proposition $\exists r,\ n = r + r$).

\subsection{Linear algebra: maps, subspaces, quotients, products}\label{ch01:sec:linalg}

\textbf{Maps.} \lean{𝒱 →ₗ[ℂ] 𝒲} is the type of $\C$-linear maps; a term is a function bundled with
the proofs that it is additive and $\C$-homogeneous. \lean{𝒱 ≃ₗ[ℂ] 𝒲} is the type of linear
isomorphisms, bundled with their inverses. \lean{Module.End ℂ 𝒱} (\lean{End ℂ 𝒱}) is
$\End_\C(\choneV)$, with \lean{*} composition, \lean{1} the identity, \lean{⁅a, b⁆} $=ab-ba$, and
\lean{D ^ N} the $N$-th power. \lean{LinearMap.id} is the identity map. A curried
\lean{M →ₗ[ℂ] M →ₗ[ℂ] ℂ} is a \emph{bilinear} map $M\times M\to\C$ (linear in each variable
separately; there is no complex conjugation anywhere), and \lean{LinearMap.BilinForm ℂ 𝒱} is the
abbreviation for \lean{𝒱 →ₗ[ℂ] 𝒱 →ₗ[ℂ] ℂ}. \lean{Module.Dual ℂ X} is the dual space
\lean{X →ₗ[ℂ] ℂ}. \lean{Basis ι ℂ X} is a basis of $X$ indexed by $\iota$.

\textbf{Subspaces as a lattice.} \lean{Submodule ℂ 𝒱} is the set of subspaces of $\choneV$, a complete
lattice under inclusion \lean{≤}: \lean{⊥} $=\{0\}$, \lean{⊤} $=\choneV$, \lean{S ⊓ T} $=S\cap T$,
\lean{S ⊔ T} $=S+T$, and \lean{⨆ i, S i} (``\lean{iSup}'') is the sum $\sum_iS_i$ of a family.
The double supremum \lean{⨆ (k : ℤ) (_ : k < n), W.grading k} in \lean{Fl} above is the sum over all
$k<n$: the second binder ranges over proofs of $k<n$, and the anonymous name \lean{_} says the proof
is not used. \lean{Submodule.span ℂ s} is the span of a set $s$, and \lean{ℂ ∙ v} the line $\C v$.
For a linear map $f$, \lean{LinearMap.ker f} and \lean{LinearMap.range f} are its kernel and image;
\lean{S.comap f} is the preimage $f^{-1}(S)$ and \lean{S.map f} the image $f(S)$; \lean{S.subtype}
is the inclusion $S\hookrightarrow\choneV$. \lean{↥S} is the subspace $S$ regarded as a vector space
(its elements are pairs of a vector and a proof of membership); the coercion is inserted
automatically.

\textbf{Quotients.} \lean{X ⧸ S} is the quotient space $X/S$. A typical construction, the cohomology
of a differential:
\noindent File \leanpath{VOA/\allowbreak Core/\allowbreak BRSTCohomology.lean}.
\begin{leancode}
def kerQ : Submodule ℂ 𝒱 := LinearMap.ker C.Q
def imQ : Submodule ℂ 𝒱 := LinearMap.range C.Q
def imQ' : Submodule ℂ C.kerQ := C.imQ.comap C.kerQ.subtype
abbrev Coh : Type u := C.kerQ ⧸ C.imQ'
\end{leancode}
Read: $\ker Q$ and $\operatorname{im}Q$ are subspaces of $\choneV$; since the quotient must be taken
inside $\ker Q$, the image is first pulled back along the inclusion $\ker Q\hookrightarrow\choneV$
(\lean{comap} along \lean{subtype}), giving $\operatorname{im}Q$ as a subspace of $\ker Q$; and the
cohomology is $\ker Q/\operatorname{im}Q$.
Besides quotient \emph{spaces}, Lean has quotient \emph{types} of a set by an equivalence relation:
\lean{Quotient.mk} sends an element to its class, \lean{Quotient.lift f h} is the function induced on
classes by an $f$ that respects the relation (\lean{h} is the proof that it does), and
\lean{Quotient.sound} says that related elements have equal classes --- it is derived from the axiom
\lean{Quot.sound} of Section~\ref{ch01:sec:trust}.

\textbf{Gradings.} A grading is a family \lean{grading : ℤ → Submodule ℂ 𝒱}, and
\lean{DirectSum.IsInternal grading} is the proposition that $\choneV = \bigoplus_nV_n$ (the natural map
from the external direct sum is bijective). \lean{DirectSum ℂ X} with \lean{X : ι → Type} is the
external direct sum $\bigoplus_iX_i$; the restricted dual of a graded module is built with it:
\noindent File \leanpath{VOA/\allowbreak Modules/\allowbreak Contragredient.lean}.
\begin{leancode}
def RestrictedDual (N : OrdVModule W M) : Type u := DirectSum ℂ (fun α => Module.Dual ℂ (N.weight α))
\end{leancode}
which is $M' = \bigoplus_\alpha M_\alpha^*$.

\textbf{Eigenspaces, exactness, complements.} For \lean{T : Module.End ℂ X},
\lean{Module.End.eigenspace T μ} is $\ker(T-\mu)$ and \lean{T.maxGenEigenspace μ} the generalised
eigenspace $\{x : (T-\mu)^kx = 0\text{ for some }k\}$; the library's logarithmic modules are graded by
the latter (Chapter~\ref{ch:L12}). \lean{IsCompl P Q} says $X = P\oplus Q$ ($P\cap Q = 0$ and
$P+Q = X$). \lean{Nontrivial X} says $X$ has two distinct elements, i.e.\ for a vector space $X\ne0$.
Exactness of a pair of module maps is the library's own predicate, stated pointwise without category
theory:
\noindent File \leanpath{VOA/\allowbreak Modules/\allowbreak OrdinaryExact.lean}.
\begin{leancode}
def ExactAt (f : Hom A B) (g : Hom B C) : Prop :=
  ∀ n : N, g.toLinearMap n = 0 ↔ ∃ m : M, f.toLinearMap m = n
\end{leancode}
i.e.\ $\ker g = \operatorname{im}f$ (Chapter~\ref{ch:L12}).

\textbf{More linear-algebra constructions.} \lean{LinearIndependent ℂ v} says the family $v$ is
linearly independent; \lean{Module.Basis} (also written \lean{Basis}) is a basis, \lean{Module.Basis.mk} builds a basis from an independent spanning family,
\lean{b.repr x} is the coordinate vector of $x$ in the basis $b$, and \lean{Pi.basisFun} is the
standard basis of $\C^n$. \lean{Module.Finite R X} says $X$ is finitely generated (for a vector space:
finite-dimensional). \lean{Submodule.Quotient.mk x} is the class of $x$ in a quotient, and
\lean{S.liftQ f h} the map induced on $X/S$ by an $f$ vanishing on $S$. \lean{LinearEquiv.ofBijective},
\lean{LinearEquiv.ofInjective} (onto the image), \lean{LinearEquiv.ofEq} (between equal subspaces) and
\lean{LinearEquiv.ofInvolutive}
build isomorphisms; \lean{LinearMap.fst} is the first projection, \lean{LinearMap.toSpanSingleton ℂ X x}
the map $c\mapsto cx$, \lean{LinearMap.smulRight f x} the map $y\mapsto f(y)\,x$, \lean{LinearMap.mk₂}
a bilinear map built from a function, \lean{LinearMap.mul} and \lean{LinearMap.mulLeft} multiplication
in an algebra, \lean{Algebra.lsmul} scalar multiplication as a linear map, \lean{B.compl₁₂ f g} the
bilinear form $(x,y)\mapsto B(fx,gy)$, \lean{B.IsSymm} and \lean{B.Nondegenerate} symmetry and
nondegeneracy of a bilinear form, \lean{Matrix.toBilin'} the bilinear form of a matrix,
\lean{Matrix.of} a matrix given by its entries, \lean{Matrix.rank} the rank and \lean{M}\lean{ᵀ} (postfix \lean{ᵀ}) the transpose.
\lean{restrictScalars} regards a $\C$-linear object as linear over a smaller ring. \lean{TensorProduct.mk},
\lean{TensorProduct.map}, \lean{TensorProduct.lift}, \lean{TensorProduct.congr}, \lean{TensorProduct.lid},
\lean{LinearMap.lTensor}, \lean{LinearMap.rTensor} are
respectively $(a,b)\mapsto a\otimes b$, $f\otimes g$, the linear map induced by a bilinear one,
$f\otimes g$ for isomorphisms, $\C\otimes M\cong M$, $\mathrm{id}\otimes f$ and $f\otimes\mathrm{id}$.
On direct sums, \lean{DirectSum.decompose} (and \lean{DirectSum.decomposeLinearEquiv}) is the isomorphism
$V\cong\bigoplus_nV_n$ of a decomposition, \lean{DirectSum.toModule} the map out of a direct sum given
on summands, \lean{DFinsupp.lsum} its finitely-supported analogue, and
\lean{DirectSum.IsInternal.chooseDecomposition} a decomposition \emph{chosen} from the bare statement
\lean{IsInternal} (Section~\ref{ch01:sec:choice}). On finitely supported functions
(Section~\ref{ch01:sec:sums}): \lean{Finsupp.single i a} is $a$ at $i$ and $0$ elsewhere (and
\lean{Pi.single} the same for ordinary functions), \lean{Finsupp.linearCombination} is
$c\mapsto\sum_ic_iv_i$, \lean{Finsupp.lsum} and \lean{Finsupp.lapply} sum and evaluation as linear
maps, \lean{Finsupp.domLCongr} re-indexing, and \lean{Finsupp.weight w} the weighted degree
$\sum_i\mu(i)w(i)$ of a multi-index.

\textbf{Dimension.} \lean{Module.finrank ℂ X} is $\dim X$ when finite and $0$ otherwise;
\lean{FiniteDimensional ℂ X} is the proposition that $X$ is finite-dimensional. A conclusion
\lean{finrank ℂ X = d} with $d\ne0$ implies finite-dimensionality; with $d = 0$ it does not.

\textbf{Products and tensor products.} \lean{A × B} is the product (direct sum) of two spaces, with
elements \lean{(a, b)} and projections \lean{x.1}, \lean{x.2}; \lean{S.prod T} is the subspace
$S\times T$. The Moonshine module is built this way:
\noindent File \leanpath{VOA/\allowbreak Moonshine/\allowbreak VnatCarrier.lean}.
\begin{leancode}
abbrev Vnat : Type := SkewVP × TP

noncomputable def gradingNat (n : ℤ) : Submodule ℂ Vnat :=
  (leechVOAPlus.grading n).prod (plusModule.weight (n : ℂ))

noncomputable def vacNat : Vnat := (leechVOAPlus.vacuum, 0)
\end{leancode}
so $\Vnat = V_\Leech^+\oplus(V_\Leech^T)^+$, graded piecewise, with vacuum $(\chonevac,0)$ (the
names \lean{SkewVP}, \lean{TP} are defined in Chapter~\ref{ch:10}). \lean{A ⊗[ℂ] B} is the tensor
product over $\C$ and \lean{a ⊗ₜ[ℂ] b} a pure tensor:
\noindent File \leanpath{VOA/\allowbreak Core/\allowbreak TensorVOA.lean}.
\begin{leancode}
def tenSub (P : Submodule ℂ A) (Q : Submodule ℂ B) : Submodule ℂ (A ⊗[ℂ] B) :=
  Submodule.span ℂ {t : A ⊗[ℂ] B | ∃ a ∈ P, ∃ b ∈ Q, a ⊗ₜ[ℂ] b = t}
\end{leancode}
which is the image of $P\otimes Q$ in $A\otimes B$, written as the span of a set given by set-builder
notation \lean{{t | ...}}; \lean{∃ a ∈ P, ...} abbreviates $\exists a,\ a\in P\wedge\dots$.

\textbf{Matrices.} \lean{Matrix ι ι ℂ} is a square matrix indexed by a finite type $\iota$, and
\lean{M.det} its determinant. A theorem combining several of these notions:
\noindent File \leanpath{VOA/\allowbreak Core/\allowbreak GradedFormMatrix.lean}.
\begin{leancode}
theorem gramMatrix_det_ne_zero_iff (hB : B.IsFullyInvariant) (hn : B.Normalized) :
    (gramMatrix B n e).det ≠ 0 ↔ maxGradedIdeal W ⊓ W.grading n = ⊥ := by
\end{leancode}
Its shape: two hypotheses on an invariant form \lean{B} (named \lean{hB} and \lean{hn}; the objects
\lean{B}, \lean{n}, \lean{e} come from \lean{variable} lines of the file), and a conclusion
``$\det\ne0$ if and only if the maximal graded ideal meets $V_n$ only in $0$''. The mathematics is in
Chapter~\ref{ch:L04}.

\subsection{Sums and finitely supported functions}\label{ch01:sec:sums}

\lean{∑ i ∈ s, f i} is the sum over a finite set \lean{s : Finset ι}; \lean{Finset.range N} is
$\{0,\dots,N-1\}$, so \lean{∑ i ∈ Finset.range N, f i} $=\sum_{i=0}^{N-1}f(i)$, and
\lean{Finset.univ} is the set of all elements of a finite type. \lean{∑ᶠ i, f i} (\lean{finsum}) is
the sum of a family indexed by an arbitrary type, defined to be the actual sum when the support
$\{i : f(i)\ne0\}$ is finite and \textbf{$0$ otherwise}. The library uses \lean{∑ᶠ} to write
identities such as the Borcherds identity as sums over all $i\in\Z$; every such use requires a
separate proof that the support is finite, and the fidelity checks point to it (Chapter~\ref{ch:L01},
warning~\ref{l01:warn:finsum}). \lean{ι →₀ R} (\lean{Finsupp}) is the type of \emph{finitely
supported} functions $\iota\to R$ --- a function together with a finite set outside which it
vanishes; \lean{μ.sum g} is the finite sum $\sum_{i}g(i,\mu(i))$ over the support:
\noindent File \leanpath{VOA/\allowbreak Modularity/\allowbreak GrowthBound.lean}.
\begin{leancode}
def multSize {g : ℕ} (μ : Fin g × ℕ →₀ ℕ) : ℕ := μ.sum (fun p m => m * (p.2 + 1))
\end{leancode}
Here $\mu$ is a multiplicity function on pairs $p = (i,m)\in\{0,\dots,g-1\}\times\N$, and the value is
$\sum_p\mu(p)\,(p_2+1)$. The anonymous function \lean{fun p m => ...} is $(p,m)\mapsto m(p_2+1)$.
Infinite sums of convergent series (\lean{∑'}, \lean{tsum}) also occur, in the analytic chapters;
\lean{∑'} of a non-summable family is again $0$ by convention; \lean{HasProd f a} is the proposition that
the infinite product of $f$ converges (unconditionally) to $a$, the honest alternative to a junk-valued
\lean{∏'}. \lean{∏ i ∈ s, f i} is a finite product.

\textbf{Supports and finiteness.} \lean{Function.support f} is the set $\{i : f(i)\ne0\}$. For a set
\lean{S : Set X}, \lean{S.Finite} and \lean{S.Infinite} are the propositions that $S$ is finite,
infinite; \lean{Set.toFinite S} is the proof of \lean{S.Finite} when $X$ itself is finite. For a type,
\lean{Finite X} says $X$ has finitely many elements and \lean{Fintype X} supplies an explicit
enumeration. \lean{Nat.card X}, \lean{Set.ncard S} and \lean{s.card} (for a \lean{Finset}) count
elements; \textbf{the first two are $0$ for infinite sets}, the same junk convention as
\lean{finrank}. \lean{Set X} is the type of subsets of $X$, \lean{Set.univ} the whole of $X$,
\lean{Set.range f} the image of $f$; \lean{∈}, \lean{∉}, \lean{⊆}, \lean{∩}, \lean{∪} have their usual
meaning. \lean{Finset.univ.filter p} is the finite set of elements satisfying $p$, and
\lean{Finset.antidiagonal n} the set of pairs $(a,b)\in\N^2$ with $a+b = n$. \lean{List X} is the type
of finite sequences, with \lean{List.replicate n x} the constant list of length $n$ and
\lean{l.flatMap f} the concatenation of the lists $f(x)$, $x\in l$; \lean{Finset.toList} lists a finite set.

\subsection{Polynomials, power series, groups}\label{ch01:sec:algebra}

\lean{Polynomial ℂ} is $\C[x]$ and \lean{MvPolynomial σ ℂ} the polynomials in variables indexed by
$\sigma$; \lean{Polynomial.aeval a P} evaluates $P$ at an element $a$ of an algebra.
\lean{PowerSeries ℂ} is $\C[[q]]$, with \lean{PowerSeries.X} the variable $q$ and
\lean{PowerSeries.coeff n f} the coefficient of $q^n$. The power series in the headline theorem is
\noindent File \leanpath{VOA/\allowbreak Modularity/\allowbreak JFunction.lean}.
\begin{leancode}
noncomputable def JShift : PowerSeries ℂ := jShift - (744 : ℂ) • PowerSeries.X
\end{leancode}
whose meaning --- $q\,J(q) = q\,(j(q)-744) = 1 + 196884\,q^2+\cdots$, with \lean{jShift} the series of
$q\,j(q)$ ---
is established in Chapter~\ref{ch:05}. A definition may need a structure on a type that is not
found automatically; it is then supplied locally with \lean{letI}:
\noindent File \leanpath{VOA/\allowbreak Modules/\allowbreak ZhuPolynomialSimplicity.lean}.
\begin{leancode}
def ZhuPolynomialRelation (P : Polynomial ℂ) : Prop :=
  letI := instRingA W
  letI := instAlgebraA W
  Polynomial.aeval (mk W W.conformalVector) P = 0
\end{leancode}
This says $P([\omega]) = 0$ in Zhu's algebra $A(W)$, where the ring and algebra structures of $A(W)$
are installed by the two \lean{letI} lines (Chapter~\ref{ch:L12}). \lean{haveI} is the same for a
proposition. Finally, \lean{Subgroup G} is the type of subgroups of a group $G$, and
\lean{Subgroup.zpowers g} the cyclic subgroup $\langle g\rangle$:
\noindent File \leanpath{VOA/\allowbreak Moonshine/\allowbreak LatticeInvolution.lean}.
\begin{leancode}
noncomputable def leechVOAPlus :
    VertexOperatorAlgebra (orbifold (Subgroup.zpowers leechTheta)).carrier :=
  orbifoldVOA (Subgroup.zpowers leechTheta)
\end{leancode}
which is $V_\Leech^+ = V_\Leech^{\langle\theta\rangle}$, the fixed points of the involution $\theta$,
as a VOA on the carrier (\lean{.carrier}) of the fixed-point subspace (Chapter~\ref{ch:04}).

\textbf{Rings, units, numbers.} \lean{CommRing R} is a commutative ring structure; \lean{algebraMap R A}
is the structure map $R\to A$ of an $R$-algebra; \lean{Subalgebra R A} is a subalgebra. \lean{A →ₐ[ℂ] B}
and \lean{A ≃ₐ[ℂ] B} are algebra homomorphisms and isomorphisms (\lean{AlgEquiv.ofAlgHom} builds one from
two mutually inverse homomorphisms), \lean{A →+ B} an additive homomorphism (\lean{f.toAddMonoidHom}
forgets the rest of a structured map), and \lean{Additive G} the additive version of a multiplicative
group. \lean{AddSubgroup G} is an additive subgroup and \lean{Equiv.neg} negation as a bijection.
\lean{IsUnit a} says $a$ is invertible, \lean{R}\lean{ˣ} (postfix \lean{ˣ}) is the group of units, \lean{a⁻¹} the inverse (with the
junk value \lean{0⁻¹ = 0} in a field), \lean{a ^ n} a power (\lean{zpow} for integer exponents), and
\lean{Commute a b} says $ab = ba$. \lean{Odd n}, \lean{a ∣ b} (divides), \lean{Nat.Coprime m n},
\lean{n.factorial} (also \lean{n !}) and \lean{Nat.choose n k} (written \lean{n.choose k}) are the usual
integer notions. \lean{Complex.I} is $i$, \lean{Complex.exp} the exponential, \lean{Real.pi} is $\pi$,
\lean{Real.sqrt} the square root (junk value $0$ on negative numbers), and \lean{CFC.sqrt} the square root
of a positive semidefinite matrix by the continuous functional calculus. \lean{f ∘ g} is composition,
\lean{∃! x, P x} unique existence, \lean{x.val} the underlying element of an element of a subtype (or of
\lean{Fin n}), \lean{Subtype.ext} says two such elements are equal when their values are, and
\lean{ofLex} strips the lexicographic-order synonym from a tuple.

\textbf{Polynomials and series.} \lean{Polynomial.X} and \lean{Polynomial.C c} are the variable and the
constant $c$; \lean{Polynomial.eval x P} is $P(x)$; \lean{P.natDegree} is the degree (junk value $0$ for
$P = 0$) and \lean{P.leadingCoeff} its leading coefficient. For several variables, \lean{MvPolynomial.X i},
\lean{MvPolynomial.C}, \lean{MvPolynomial.aeval}, \lean{MvPolynomial.lcoeff m} (the coefficient of the
monomial $m$, as a linear map; \lean{lcoeff} also exists for \lean{Polynomial} and \lean{PowerSeries}),
\lean{MvPolynomial.pderiv i} ($\partial/\partial x_i$) and
\lean{MvPolynomial.weightedHomogeneousSubmodule R w d} (the polynomials of weighted degree $d$).
\lean{PowerSeries.mk a} is the series $\sum a_nq^n$. Mathlib's modular forms appear as
\lean{ModularForm Γ k} (weight-$k$ forms for a group $\Gamma$) and their $q$-expansions
\lean{qExpansion}; the Eisenstein-series lemma \lean{EisensteinSeries.E_qExpansion_coeff} gives their
coefficients (Chapters~\ref{ch:L14} and~\ref{ch:05}).

\textbf{Lie theory.} \lean{LieRing L} and \lean{LieAlgebra ℂ L} give a Lie bracket \lean{⁅x, y⁆};
\lean{L →ₗ⁅ℂ⁆ L'} (\lean{LieHom}) is a Lie algebra homomorphism, \lean{LieSubalgebra ℂ L} a Lie
subalgebra, \lean{UniversalEnvelopingAlgebra ℂ L} is $U(L)$, \lean{IsSl2Triple h e f} says $(h,e,f)$
satisfies the $\mathfrak{sl}_2$ relations, and \lean{oneCochain}, \lean{twoCochain}, \lean{twoCocycle},
\lean{LieAlgebra.ofTwoCocycle} are Mathlib's Lie algebra cochains, cocycles and the central extension
defined by a cocycle --- how the Virasoro algebra is built from the Witt algebra (Chapter~\ref{ch:L02}).
\lean{IsSimpleModule R M} and \lean{IsSemisimpleRing R} are the module-theoretic notions; for a lattice
of subspaces, \lean{sSup} and \lean{sInf} are the sum and intersection of a set of subspaces.
\lean{Semiring}, \lean{AddCommMonoid}, \lean{LinearOrder} and \lean{NoZeroDivisors} are the evident
algebraic and order classes, appearing as instance arguments. \lean{Ideal R} is the type of ideals of a
ring, \lean{Ideal.span s} the ideal generated by $s$, \lean{RingHom.ker f} the kernel of a ring
homomorphism, and \lean{I.mkQ} (for a subspace or ideal $I$) the quotient map $X\to X/I$. For
polynomials, \lean{P.derivative} is $P'$ and \lean{P.eval x} evaluation, so \lean{M.det.eval x} and
\lean{M.det.derivative.eval x} evaluate a polynomial determinant and its derivative; \lean{MvPolynomial.weightedHomogeneousComponent w d P} is the weighted
degree-$d$ part of $P$. \lean{Unit} is the one-element type, \lean{l.sort r} sorts a list or finite set
by a relation $r$, and \lean{cmp a b} compares two elements of a linear order. Hypothesis names that are
words (\lean{hsum}, \lean{hne}, \lean{hB}) are just labels chosen by the authors; only their types
matter. \lean{Submodule.Quotient.restrictScalarsEquiv} and lemma names such as
\lean{Submodule.mem_span_singleton_self} ($x\in\C x$) occur in a few statements and mean what they say.

\subsection{How a Lean statement can fail to mean what it says}\label{ch01:sec:failures}

Everything above is vocabulary. The point of learning it is to recognise the ways in which a
correct, machine-checked statement can still fail to be the intended theorem. The book's fidelity
checks are organised around the following list; each entry has occurred in this library, and the
instances are catalogued in Chapter~\ref{ch:15}, Sections~\ref{ch15:sec:deadends}--\ref{ch15:sec:debt}.
\begin{enumerate}[label=(F\arabic*)]
\item \textbf{Unsatisfiable structure.} Its axioms are contradictory, so every theorem about it is
  vacuous. Instances: the original \lean{TwistedModule} (repaired; Chapter~\ref{ch:06}), the
  unrepaired \lean{SuperTwistedModule.jacobi} (Chapter~\ref{ch:L17}; Section~\ref{ch15:sec:super}), and
  \lean{VOA.Unitary}, whose ``positive-definiteness'' is required of a \emph{bilinear} (not
  Hermitian) form and so cannot hold in dimension $\ge2$ (Chapter~\ref{ch:L12},
  Proposition~\ref{l12:prop:unitary}). The \lean{¬ Nonempty} theorem quoted in
  Section~\ref{ch01:sec:choice} is the library recording such a case itself.
\item \textbf{Predicate satisfied too easily.} \lean{IsLowestWeight N h} does not require the
  lowest weight space to be nonzero, so the zero module is lowest-weight of every weight
  (Chapter~\ref{ch:L12}); \lean{SelfContragredient} holds for every module (Chapter~\ref{ch:L12}).
\item \textbf{Predicate satisfied too rarely.} \lean{IsCommutative} forces $V = \C\chonevac$, so the
  ``commutativity theorems'' are statements about the trivial VOA only (Chapter~\ref{ch:L03}).
\item \textbf{Definition differs from its name.} The library's structure transcribing Gannon's axioms lacks
  translation covariance (Chapter~\ref{ch:L03}); it was called \lean{VertexAlgebra} until the release, when it was renamed
  \lean{PreVertexAlgebra} and the name \lean{VertexAlgebra} was given to the translation-covariant notion. Similarly the
  universal affine VOA $V^k(\mathfrak g)$, formerly \lean{simpleAffineVOA}, is now \lean{universalAffineVOA}
  (Chapter~\ref{ch:L23}). Still open: the super structure \lean{SuperVertexAlgebra}, which has the same omission, keeps the old-style
  name (Chapter~\ref{ch:L17}); the library's \lean{generated} subalgebra always contains $\omega$
  (Chapter~\ref{ch:L04}); its \lean{fusionCoeff} is identically $0$, a placeholder and not a fusion rule (Chapter~\ref{ch:L12}).
\item \textbf{Junk values} in a conclusion (\lean{finrank}, \lean{∑ᶠ}, truncated subtraction),
  Section~\ref{ch01:sec:junk}.
\item \textbf{A hypothesis that is never established}, making the theorem conditional although its
  axiom list is clean (Section~\ref{ch01:sec:uncond}); or a hypothesis that is redundant or
  tautological, which is harmless but misleading (for instance the hypothesis \lean{homega} still carried by the definition \lean{inner}, Remark~\ref{l04:rem:homega}).
\item \textbf{Quantifier order or a missing quantifier}: a bound that should be uniform but depends on
  the data, or the reverse (Section~\ref{ch01:sec:lean}).
\item \textbf{A conclusion satisfied by a trivial or wrong object}: the zero VOA satisfies every
  axiom; a non-simple square-zero extension has the same character as $\Vnat$
  (Chapter~\ref{ch:10}). Headline theorems therefore pin the data and assert simplicity.
\item \textbf{Choices and instances} that change meaning: a theorem about a chosen witness
  (Section~\ref{ch01:sec:choice}), or an operation whose meaning comes from an unexpected instance
  (Example~\ref{ch01:ex:instOne}, Section~\ref{ch01:sec:synonyms}).
\end{enumerate}

\subsection{A second structure: super vertex algebras}\label{ch01:sec:super}

The same reading technique applies to every structure of the library. As a second example, and
because it shows two further pieces of vocabulary (gradings indexed by $\C$ and by \lean{ZMod 2}),
here is the library's super vertex algebra.
\noindent File \leanpath{VOA/\allowbreak Core/\allowbreak SuperVertexAlgebra.lean}.
\begin{leancode}
structure SuperVertexAlgebra (𝒱 : Type u) [AddCommGroup 𝒱] [Module ℂ 𝒱] where
  grading : ℂ → Submodule ℂ 𝒱
  isInternal : DirectSum.IsInternal grading
  finrank : ∀ h, FiniteDimensional ℂ (grading h)
  parity : ZMod 2 → Submodule ℂ 𝒱
  parity_isInternal : DirectSum.IsInternal parity
  Y : 𝒱 →ₗ[ℂ] (End ℂ 𝒱)⟦z^±⟧
  vacuum : 𝒱
  vacuum_mem : vacuum ∈ grading 0
  vacuum_parity : vacuum ∈ parity 0
  grading_mode : ∀ {h k : ℂ} (u : 𝒱), u ∈ grading h → ∀ (v : 𝒱), v ∈ grading k →
    ∀ (n : ℤ), mode (Y u) n v ∈ grading (h + k - n - 1)
  parity_mode : ∀ {p q : ZMod 2} (u : 𝒱), u ∈ parity p → ∀ (v : 𝒱), v ∈ parity q →
    ∀ (n : ℤ), mode (Y u) n v ∈ parity (p + q)
  vacuum_Y : Y vacuum = (1 : (End ℂ 𝒱)⟦z^±⟧)
  creation : ∀ u : 𝒱, mode (Y u) (-1) vacuum = u ∧ ∀ n : ℤ, 0 ≤ n → mode (Y u) n vacuum = 0
  locality : ∀ {p q : ZMod 2} (u : 𝒱), u ∈ parity p → ∀ (v : 𝒱), v ∈ parity q →
    Localε (koszul p q) (Y u) (Y v)
  regularity : ∀ u v : 𝒱, ∃ N : ℤ, ∀ n : ℤ, N ≤ n → mode (Y u) n v = 0
\end{leancode}

\begin{dictionary}
\textbf{How to read it.} Data: two gradings (\lean{grading}, indexed by $\C$ so that half-integer
weights are allowed while modes remain indexed by $\Z$; \lean{parity}, indexed by \lean{ZMod 2}), the
vertex operator \lean{Y} and the \lean{vacuum}. Everything else is an axiom. Note the field
\lean{finrank}: despite its name it is the \emph{proposition} that each graded piece is
finite-dimensional, not a number. \lean{creation} is a conjunction \lean{∧} of two clauses.
\lean{locality} is super-locality with the Koszul sign \lean{koszul p q} $=(-1)^{pq}$, defined in
Chapter~\ref{ch:L01} (\lean{Localε}, Definition~\ref{l01:def:bracketeps}); \lean{regularity} has the
quantifier order ``for all $u,v$ there is $N$''. No relation between weight and parity
(spin--statistics) is imposed; the theory of super vertex algebras is Chapter~\ref{ch:L17}. The
Moonshine construction does not use this structure.
\end{dictionary}

\begin{physics}
A super vertex algebra is the chiral algebra of a theory with fermions: odd states create fermionic
fields, which anticommute at spacelike separation, whence the sign $(-1)^{pq}$ in locality. Free
fermions and the $bc$ and $\beta\gamma$ ghost systems of string theory are examples.
\end{physics}

\begin{status}
A definition. \texttt{\#print axioms VOA.SuperVertexAlgebra} lists
\lean{[propext, Classical.choice, Quot.sound]} (run for this chapter).
\end{status}

\subsection{Reading the headline statement}\label{ch01:sec:headline}

Here is the main theorem of the Moonshine part of the library, as it stands in the source.
\noindent File \leanpath{VOA/\allowbreak Moonshine/\allowbreak TopPairs.lean}.
\begin{leancode}
theorem moonshine :
    ∃ V : VertexOperatorAlgebra Vnat,
      V.Y = Ynat VnatAxioms.wbarL ∧ V.vacuum = vacNat ∧ V.conformalVector = omegaNat ∧
      V.centralCharge = 24 ∧
      V.grading 1 = ⊥ ∧
      Module.finrank ℂ (V.grading 2) = 196884 ∧
      (∀ n : ℕ, (Module.finrank ℂ (V.grading (n : ℤ)) : ℂ) = PowerSeries.coeff n JShift) ∧
      V.IsSimple :=
  moonshine_of_topPairs (by norm_num) topPairs_one
\end{leancode}
Its full name is \lean{VOA.Moonshine.moonshine}. It has \emph{no arguments}: there is nothing
before the colon. Its statement is an existential, \lean{∃ V : VertexOperatorAlgebra Vnat, ...},
whose body is a conjunction of eight clauses. Read it as: \emph{there is a vertex operator algebra
structure $V$ on the space \lean{Vnat} such that}
\begin{enumerate}[label=(\arabic*)]
\item its vertex operator is the explicitly constructed block map \lean{Ynat VnatAxioms.wbarL};
\item its vacuum is the explicit vector \lean{vacNat};
\item its conformal vector is the explicit vector \lean{omegaNat};
\item its central charge is $24$;
\item its weight-one space is zero, $V_1 = 0$;
\item $\dim V_2 = 196884$;
\item for every $n\in\N$, $\dim V_n$ (cast to $\C$) is the $q^n$-coefficient of the power series
  \lean{JShift};
\item it is simple.
\end{enumerate}
The term after \lean{:=} is the whole proof: it applies an earlier theorem,
\lean{moonshine_of_topPairs}, to a numerical fact checked by \lean{norm_num} and to the theorem
\lean{topPairs_one}. Both are explained in Chapter~\ref{ch:14}.
\begin{dictionary}
\textbf{What to notice.} (a) The first three clauses \emph{pin} the structure to the explicit
Frenkel--Lepowsky--Meurman data, so the existential cannot be satisfied by some unrelated object
on the same space. (b) Clause (7) contains a \emph{cast}: \lean{Module.finrank} is a natural
number, and \lean{(... : ℂ)} turns it into a complex number so that it can be compared with a
coefficient of a power series over $\C$; and \lean{(n : ℤ)} casts $n\in\N$ to an integer weight.
(c) The source writes \lean{Module.finrank ℂ (V.grading 2)}, but \lean{V.grading 2} is a subspace,
not a type: Lean silently inserts the coercion \lean{↥}, and Lean's own printout of the statement
shows \lean{Module.finrank ℂ ↥(V.grading 2)}. (d) Clause (8) is essential: a non-simple VOA (the
``square-zero extension'', Chapter~\ref{ch:10}) has exactly the same central charge and graded
dimension. (e) Each of the defined terms \lean{Vnat}, \lean{Ynat}, \lean{vacNat},
\lean{omegaNat}, \lean{JShift}, \lean{grading} and \lean{IsSimple} must be traced to its
definition. The places are listed in Table~\ref{ch01:tab:notions}, and the complete clause-by-clause
fidelity check is Chapter~\ref{ch:14}, Section~\ref{ch14:sec:decode}.
\end{dictionary}

\begin{physics}
The theorem says that the chiral conformal field theory of the $\Z/2$-orbifold of 24 free bosons on
the Leech torus exists as a mathematical object, has central charge $24$, has no conserved weight-one
currents ($V_1 = 0$: no continuous symmetry), has $196884$ states of weight two, and has partition
function $J(q)$; simplicity says that the theory has a unique vacuum sector with no decoupled
subsector.
\end{physics}

\begin{status}
Unconditional: the theorem has no hypotheses. \texttt{\#print axioms VOA.Moonshine.moonshine}, run for
this chapter, gives
\begin{plaincode}
'VOA.Moonshine.moonshine' depends on axioms: [propext, Classical.choice, Quot.sound]
\end{plaincode}
so the proof uses neither \lean{sorry} nor any project-specific axiom (the library has none;
Chapter~\ref{ch:L02}).
\end{status}

\subsection{Namespaces and full names}

Declarations live in namespaces, and the full name records them:
\lean{VOA.VertexOperatorAlgebra.L} is the definition \lean{L} in the namespace
\lean{VertexOperatorAlgebra} inside \lean{VOA}. The whole library is inside \lean{VOA}; the Moonshine
part is mostly inside \lean{VOA.Moonshine}. Inside a namespace, and after an \lean{open} command,
Lean lets one write short names, so a printed statement may say \lean{finrank} for
\lean{Module.finrank} or \lean{End} for \lean{Module.End}. The book gives full names in the text and
the file in which each declaration lives.

%======================================================================
\section{What the machine guarantees}\label{ch01:sec:trust}
%======================================================================

\subsection{The kernel}

When Lean accepts a file, its \emph{kernel} --- a small program that implements the rules of the type
theory --- has checked every proof term in it. Everything a user writes, including the
automation (``tactics'') that produces proofs, is outside the kernel and does not need to be
trusted: if a tactic produced a wrong proof, the kernel would reject it. What a reader must trust is
therefore (i) the kernel and the logic, (ii) the axioms, and (iii) the \emph{statements}. This book
is about (iii).

\subsection{Escape hatches: \texttt{sorry} and friends}

Lean has an escape hatch, \lean{sorry}, which closes any goal without proof; internally it is an
axiom \lean{sorryAx}. A statement proved with \lean{sorry} is not proved. The project's policy
forbids \lean{sorry}, \lean{admit}, \lean{native_decide} (which trusts the compiler rather than the
kernel), \lean{implemented_by} and \lean{stop} in library code; the repository checker enforces it.

\subsection{Axioms and \texttt{\#print axioms}}

Lean's logic uses three standard axioms beyond its type theory: \lean{propext} (logically equivalent
propositions are equal), \lean{Quot.sound} (quotient types behave as quotients) and
\lean{Classical.choice} (the axiom of choice, from which the law of the excluded middle follows).
All of classical mathematics and all of Mathlib use them. A user may add further axioms with the
keyword \lean{axiom}, and Lean will then accept anything that follows from them; a false axiom makes
everything provable. The command \texttt{\#print axioms foo} traverses the proof of \lean{foo} down
to the kernel and lists every axiom it uses.

\textbf{The library contains no project axiom.} Until 19 October 2026 it
contained exactly one, \lean{VOA.pbw}; that declaration is now a theorem with the same name and type.
\noindent File \leanpath{VOA/\allowbreak Foundations/\allowbreak AffineLie.lean}.
\begin{leancode}
theorem pbw : PBWPrinciple :=
  ⟨fun κ hκ k => AffineLie.vacuumVector_ne_zero_of_induction κ hκ k,
    fun p rho => PBWTheorem.bottom_injective p rho⟩
\end{leancode}
\lean{PBWPrinciple} bundles two consequences of the Poincar\'e--Birkhoff--Witt theorem: that the
vacuum vector of an affine vacuum module is nonzero, and that a module embeds in the module induced
from it. Mathlib does not contain PBW, and the library proves the needed relative form itself
(Chapter~\ref{ch:L02}, Section~\ref{l02:sec:pbwproof}). Its exact statement, its proof, and which
parts of the library use it are in Chapter~\ref{ch:L02}, Section~\ref{l02:sec:axiom}. Earlier
editions flagged every result that depended on it; those flags are historical
(Section~\ref{l02:sec:history}). For example, run for this book at an earlier commit, after \lean{VOA.pbw} became a theorem (not re-run for this edition):
\begin{plaincode}
'VOA.AffineLie.vacuumVector_ne_zero' depends on axioms: [propext, Classical.choice, Quot.sound]
\end{plaincode}
The construction of $\Vnat$ does not use it, as the axiom report of the headline theorem
(Section~\ref{ch01:sec:headline}) certifies.

\subsection{A clean axiom list does not make a theorem unconditional}\label{ch01:sec:uncond}

\texttt{\#print axioms} reports \emph{axioms}, not \emph{hypotheses}. A theorem
\lean{theorem bar (h : SomeHypothesis) : Conclusion} can have a perfectly clean axiom list and still
prove nothing unconditional: \lean{h} is an explicit argument, and the theorem says only ``if
\lean{SomeHypothesis} then \lean{Conclusion}''. The book therefore uses two labels, and never
blurs them:
\begin{itemize}
\item \textbf{unconditional}: the statement has no hypotheses beyond the objects it speaks about and
  their defining structure, and its axioms are the standard three;
\item \textbf{\texttt{\#print axioms}-clean, conditional on $H$}: the statement has an explicit
  hypothesis $H$ that is not established anywhere, so it proves $H\Rightarrow C$, not $C$;
\end{itemize}
(Earlier editions had a third label, ``uses \lean{VOA.pbw}'', for results whose axiom list contained
the then project axiom; it no longer applies.)
Hypotheses that are part of the mathematical statement (``let $V$ be simple'', ``let
$L_1V_1 = 0$'') are of course fine; the question is always whether a hypothesis is an intended
condition of the theorem or an unproved input. The library contains many results of the second kind
--- for example Zhu-type modular invariance under Etingof-finiteness in
Chapter~\ref{ch:L16} takes hypotheses \lean{hfib} and \lean{hrank} that are discharged only for the trivial
VOA --- and the status boxes name them. (Such inputs are sometimes later proved: the Arakawa--Kawasetsu weight grading
\lean{AKWeightGrading W}, long carried as such a hypothesis, is now the theorem \lean{VOA.akWeightGrading} for every VOA.) The headline theorem \lean{VOA.Moonshine.moonshine} has no hypotheses at all.

\subsection{The aggregate build}

\texttt{lake} is Lean's build tool. \texttt{lake build} elaborates every file of the library in
dependency order and has the kernel check every proof; ``green'' means it finished with no errors.
The project's rule is that a result counts only when the \emph{aggregate} build of the umbrella
module \lean{VOA} is green: a file that compiles in isolation might not be imported by the library
at all. The repository checker \texttt{scripts/check\_repository.py} verifies that every library
file is reachable from the umbrella, that no forbidden escape hatch occurs, and that no project
\lean{axiom} declaration exists.

\begin{warning}
\textbf{Audit trap.} Running a bare \texttt{lake env lean} on a file can silently drop the option
\lean{maxSynthPendingDepth} that the library's build sets; elaboration may then fail in a way that
inserts \lean{sorryAx}. Axiom audits are therefore run only through \texttt{lake build} or through
\texttt{lake env lean -DmaxSynthPendingDepth=3}. All axiom reports in this chapter were obtained the
second way.
\end{warning}

\subsection{The book's own machine checks}\label{ch01:sec:bookchecks}

The book itself is checked by machine in three ways, by \texttt{python3 scripts/textbook/sync.py
check}.
\begin{itemize}
\item \textbf{STALE QUOTE.} Every \lean{leancode} block of every chapter is compared with the Lean
  sources (after removing comments and normalising whitespace); a block that does not occur
  verbatim is reported. This is what guarantees that the Lean code printed in the book is the code
  in the library. Tool output, such as \texttt{\#print axioms} reports, is printed in a different
  (plain) style and is not compared.
\item \textbf{MISSING PATH.} Every quoted declaration must be accompanied by a citation of the file
  that defines it (the ``File \leanpath{...}'' lines); \texttt{scripts/textbook/cite\_paths.py}
  checks this.
\item \textbf{NO PHYSICS.} Every definition and theorem must be followed by a physics box.
\end{itemize}
In addition, \texttt{sync.py} records for every \lean{.lean} file the chapter that documents it and
its hash at the documented commit, so that when the library changes it lists the new, changed and
deleted files chapter by chapter. Where a chapter needed a Lean computation of its own --- a
counterexample, or a proof that a definition is degenerate --- the computation was kept outside the
library and compiled against the library; these checks are not part of this release, and the
chapters say so where they use one. Finally, every defect of the library
found while writing or refereeing the book (a definition that is vacuous, weaker or stronger than its
name, a redundant hypothesis, a wrong docstring) is reported, with its evidence and a severity, in the
chapter where it occurs and collected in the catalogue of Chapter~\ref{ch:15}
(Sections~\ref{ch15:sec:deadends}, \ref{ch15:sec:deadends-oct} and~\ref{ch15:sec:debt}).

\subsection{The verification record for this chapter}

\emph{Run for this chapter} at an earlier commit (the snapshot against which this chapter was
first written; this record is historical and was not re-run for this edition), with
\texttt{lake env lean -DmaxSynthPendingDepth=3} on a scratch file importing the relevant modules
(the library's compiled files were present; no build was run):
\texttt{\#print axioms} of \lean{VOA.Moonshine.moonshine}, \lean{VOA.Aut},
\lean{VOA.VertexOperatorAlgebra} and \lean{VOA.SuperVertexAlgebra} each gave
\lean{[propext, Classical.choice, Quot.sound]}; \lean{VOA.mode} and \lean{VOA.FormalDist.instOne} gave no
axioms; \lean{VOA.AffineLie.vacuumVector_ne_zero} gave the standard three and \lean{VOA.pbw} (then an axiom). The
repository checker reported, in that run,
\begin{plaincode}
Repository check passed: 1867 library files scanned; 1867 reachable through 4676 local import edges; 1 sanctioned axiom declaration(s).
\end{plaincode}
In a later run, after \lean{VOA.pbw} became a theorem (also not re-run for this edition),
\lean{VOA.AffineLie.vacuumVector_ne_zero} gave the standard three axioms only, and the checker reported
``0 sanctioned axiom declaration(s)''.
\emph{Reported by the project's build record for an earlier commit and not re-run for this chapter:}
the aggregate \texttt{lake build} was green and the axiom-audit script passed. Since then the aggregate build \texttt{bash scripts/campaign\_build.sh VOA} (the full \texttt{lake build VOA} under the project's memory-capped wrapper) was green in two pre-release builds on 2026-10-28 (18{,}901 jobs, and 18{,}907 jobs after merging the then-current \texttt{origin/main}); \texttt{scripts/check\_repository.py} passed (10{,}431 library files, all reachable, 0 project axioms) and \texttt{scripts/audit/run\_axiom\_audits.py} passed (5{,}958 reports across 24 probes). The checks at the release commit are recorded in Chapter~\ref{ch:15}, Section~\ref{ch15:sec:record}, which collects the
verification record of the whole book.

\subsection{Traps that concern the whole book}\label{ch01:sec:traps}

\begin{warning}
\textbf{Recorded traps.} (1) A structure with unsatisfiable axioms makes every theorem about it
vacuous; the library's \lean{TwistedModule} had this defect and was repaired, and the same
over-quantification remains, unrepaired, in \lean{SuperTwistedModule.jacobi}
(Chapters~\ref{ch:06} and~\ref{ch:L17}). (2) Graded dimension does not characterise $\Vnat$: a
non-simple square-zero extension has the same central charge and character, so simplicity is
essential (Chapter~\ref{ch:10}). (3) \lean{Module.finrank} of an infinite-dimensional space is $0$,
and \lean{∑ᶠ} of an infinitely supported family is $0$. (4) The library's \lean{fusionCoeff} is not
the fusion rule; it is identically zero (Chapter~\ref{ch:L12}). (5) A bare \lean{PreVertexAlgebra} does
not satisfy the Borcherds identity (above, and Chapter~\ref{ch:L03}). (6) Audit only through
\texttt{lake build} or with \texttt{-DmaxSynthPendingDepth=3}.
\end{warning}

%======================================================================
\section{Global conventions}\label{ch01:sec:conventions}
%======================================================================

The following conventions are used throughout the library and the book. Each is defined in full,
with its Lean source, in the place indicated, usually with a worked example (the mode convention is
worked out in Example~\ref{ch01:ex:mode}, the odd-integer encoding in Section~\ref{ch01:sec:junk});
the table is a quick reference, not a definition. To decode an entry, read the ``Meaning'' column
as the mathematical expression that the Lean expression in the first column denotes.

\begingroup\small
\begin{longtable}{>{\raggedright\arraybackslash}p{0.27\textwidth}>{\raggedright\arraybackslash}p{0.47\textwidth}>{\raggedright\arraybackslash}p{0.18\textwidth}}
\toprule
Lean & Meaning & Defined in\\
\midrule\endhead
\lean{FormalDist 𝒲}, \lean{𝒲⟦z^±⟧} & functions $\Z\to\mathcal W$, i.e.\ series $\sum_k f(k)z^k$; \lean{f k} is the coefficient of $z^k$ & Def.~\ref{l01:def:formaldist}\\
\lean{mode f n} & $f_{(n)} = f(-n-1)$, the coefficient of $z^{-n-1}$ & Def.~\ref{l01:def:mode}\\
\lean{mode (Y u) n} & the operator $u_{(n)}$ in $Y(u,z) = \sum_n u_{(n)}z^{-n-1}$; \lean{Y u k} is $u_{(-k-1)}$ & Def.~\ref{l01:def:mode}, \ref{l03:def:VA}\\
\lean{V.smul u n v} & the $n$-th product $u_{(n)}v$ (not scalar multiplication) & Def.~\ref{l03:def:VA}\\
\lean{W.L n} & $L_n = \omega_{(n+1)}$, i.e.\ \lean{mode (Y conformalVector) (n + 1)}; $L_0 = \omega_{(1)}$, $L_{-1} = \omega_{(0)}$ & Def.~\ref{l03:def:VOA}\\
\lean{V.T} & $Tu = u_{(-2)}\chonevac$; equals $L_{-1}$ in a VOA & Def.~\ref{l03:def:smulT}\\
\lean{deriv} & formal $\partial_z$: coefficient of $z^n$ is $(n+1)f(n+1)$; in modes $(\partial f)_{(m)} = -m f_{(m-1)}$ & Def.~\ref{l01:def:mode}\\
\lean{1 : (End ℂ 𝒱)⟦z^±⟧} & the identity field (only the $z^0$ coefficient is nonzero) & Example~\ref{ch01:ex:instOne}\\
\lean{grading : ℤ → Submodule ℂ 𝒱} with \lean{DirectSum.IsInternal} & a $\Z$-grading $V = \bigoplus_n V_n$ (super: indexed by $\C$, parity by \lean{ZMod 2}) & Def.~\ref{l03:def:VA}; §\ref{ch01:sec:super}\\
\lean{N.weight α} & the weight-$\alpha$ space $M_\alpha$ of a module, $\alpha\in\C$ & Chapter~\ref{ch:L12}, §\ref{l12:sec:recall}\\
\lean{modeT r i k} & twisted-sector mode $h_i(k/2)$: half-integer modes indexed by odd integers $k\mapsto k/2$ & Chapter~\ref{ch:06}, §\ref{ch06:sec:heis}\\
\lean{TwistedModule.YM u k} & the twisted mode $u_{(k/N)}$ for an automorphism of order $N$ & Def.~\ref{ch06:def:tm}\\
\lean{gBinom n k} & $\binom nk = n(n-1)\cdots(n-k+1)/k!$ for $n\in\Z$, $k\in\N$ & Def.~\ref{l01:def:gbinom}\\
\lean{cbz n k} & $\binom nk$ with $k\in\Z$, defined to be $0$ for $k<0$ & Def.~\ref{l01:def:gbinom}\\
\lean{negOnePow n} & $(-1)^n$ & Def.~\ref{l01:def:gbinom}\\
binomial expansion & $(z_1+z_2)^n = \sum_{k\ge0}\binom nk z_1^{n-k}z_2^k$, always in nonnegative powers of the \emph{second} variable & §\ref{l01:sec:binomial}\\
\lean{Local a b} & $(z_1-z_2)^N[a(z_1),b(z_2)] = 0$ for some $N$ & Def.~\ref{l01:def:bracket}\\
\lean{Module.finrank} & dimension; $0$ for infinite-dimensional spaces & §\ref{ch01:sec:lean}\\
\lean{∑ᶠ} & finitely supported sum; $0$ if the support is infinite & warning~\ref{l01:warn:finsum}\\
\lean{ℕ}, casts & $\N\ni0$; \lean{(n : ℂ)}, \lean{↑n} cast integers to $\C$ & §\ref{ch01:sec:lean}\\
Virasoro & $[L_m,L_n] = (m-n)L_{m+n} + \frac{m^3-m}{12}c\,\delta_{m+n,0}$ & Def.~\ref{l02:def:vir}, \ref{l03:def:VOA}\\
\bottomrule
\end{longtable}
\endgroup

%======================================================================
\section{How to read the boxes}\label{ch01:sec:boxes}
%======================================================================

Four kinds of coloured box recur. They have fixed jobs.

\begin{dictionary}
\textbf{The fidelity check (grey).} This is the heart of the book. It decodes every symbol and
convention of the Lean statement; lists \emph{every} hypothesis, including those inherited from
\lean{variable} lines, and says what each means and whether it is intended, harmless, or an open
input; traces every defined term to its definition; says whether the Lean statement is the same as,
stronger than, or weaker than the textbook statement; and asks how it could have been wrong or
vacuous. A referee who disagrees with a fidelity check has found something important.
\end{dictionary}

\begin{physics}
\textbf{The physics box (blue).} Says what the statement means in two-dimensional conformal field
theory, string theory or statistical mechanics: which states, fields, operator product expansions,
symmetries, partition functions or orbifolds it is about, and why a physicist would expect it.
Heuristics are labelled as such; the physical picture never substitutes for the proof.
\end{physics}

\begin{status}
\textbf{The status box (green).} The machine status, honestly: unconditional, or conditional on named
hypotheses (listed); the result of \texttt{\#print axioms} where it was run, and whether it was run
by the chapter's author or is reported from elsewhere; (in earlier editions) whether \lean{VOA.pbw} was used; numerical
evidence clearly labelled as evidence.
\end{status}

\begin{warning}
\textbf{The warning box (red).} Traps, dead ends, corrected mistakes and defects of the library, with
the counterexample or other evidence. Every defect recorded in a warning box is also collected in
the catalogue of Chapter~\ref{ch:15} (Section~\ref{ch15:sec:deadends}).
\end{warning}

Each chapter ends with a \textbf{Referee's checklist}: a table of the chapter's headline
declarations with, for each, the intended meaning in one sentence, the hypotheses,
conditional/unconditional, and a pointer to its fidelity check. A referee who wants to audit a
chapter quickly can start from that table.

%======================================================================
\section{A map of the book}\label{ch01:sec:map}
%======================================================================

\subsection{The parts}

The book has eight parts. Parts~I--V document the general library; Part~VI formalizes recent research
theorems, and conditional reductions of further ones, on top of it; Part~VII constructs the Moonshine
module; Part~VIII records how all of it was verified.
\begin{description}
\item[Part I, Foundations.] This chapter; Chapter~\ref{ch:L01}, formal calculus and locality
  (formal distributions, modes, the delta function, generalised binomials, locality and its
  delta-function characterisation); Chapter~\ref{ch:L02}, the Lie-theoretic inputs (Witt, Virasoro
  and affine Lie algebras, vacuum modules, root data, lattices) and the PBW theorem
  \lean{VOA.pbw} (formerly a project axiom, now a theorem); Chapter~\ref{ch:L03}, the axioms of vertex algebras and VOAs, Dong's lemma,
  Goddard uniqueness, the Borcherds identity, skew-symmetry and the reconstruction theorem;
  Chapter~\ref{ch:L04}, structure theory: the Lie algebra $V_1$, automorphisms, ideals, simplicity
  and quotients, homomorphisms, generated subalgebras, invariant bilinear forms and Li's criterion
  ``simple $\iff$ nondegenerate''.
\item[Part II, Examples, deformations and rigidity.] Chapters~\ref{ch:L05}--\ref{ch:L11}: Heisenberg
  and lattice VOAs; the Virasoro VOA, its Verma modules and the Kac determinant; first-order
  deformations and the rigidity of Virasoro minimal models; affine $\mathfrak{sl}_2$ and its rigidity
  at every admissible level; filtrations, the $C_2$ algebra and presentations. A chapter on
  affine $\mathfrak{sl}_3$ is in preparation. Chapter~\ref{ch:L11} also contains a second criterion for the surjectivity of Li's
  jet map, \emph{generator-word spanning} (Section~\ref{l11:sec:genwords}): if Li's filtration equals the
  filtration by generator words, the jet map is onto, for any VOA and any finite family of homogeneous
  generators. Unlike the earlier criterion, it does not force $C_2$-cofiniteness, and it is the
  surjectivity half of the classical-freeness results of Chapters~\ref{ch:L20} and~\ref{ch:L24}.
\item[Part III, Representation theory.] Chapter~\ref{ch:L12}: weak, admissible, logarithmic and
  ordinary modules, Zhu's algebra, $C_2$-cofiniteness, contragredients, intertwining operators and
  the rationality hierarchy; Chapter~\ref{ch:L13}: Zhu's bijection via the Frenkel--Zhu tower.
\item[Part IV, Modularity.] Chapters~\ref{ch:L14}--\ref{ch:L16}: modular forms and the analytic
  toolkit; Zhu's modularity theorem; quasi-lisse vertex algebras. Chapter~\ref{ch:L16} should be read
  with its own scope: ``quasi-lisse'' is not defined in the library (the theorems assume the weaker
  finiteness of zeroth Poisson homology, \lean{EtingofFinite}). The Arakawa--Kawasetsu modular linear
  differential equation (Theorem~5.1) is proved conditionally only on \lean{EtingofFinite}: one monic MLDE, of some
  unbounded order, satisfied by every abstract analytic genus-one block at a square-bracket primary insertion
  (\lean{VOA.IsSqPrimary.exists_monicMLDEOn_all_analyticBlocks_of_etingofFinite_uncond}; its former inputs, the weight
  grading \lean{VOA.akWeightGrading} and the Serre closure \lean{VOA.e46SerreClosed}, are theorems). Unconditionally,
  $L_{-4/3}(\mathfrak{sl}_2)$ is not $C_2$-cofinite but has finite-dimensional zeroth Poisson homology, and its vacuum
  one-point blocks satisfy one monic MLDE of unspecified order (\lean{VOA.exists_monicMLDEOn_all_analyticBlocks_simpleSl2QL};
  Arakawa--Kawasetsu's order~2 is not proved).
\item[Part V, Super, logarithmic and unitary theories.] Chapters~\ref{ch:L17}--\ref{ch:L19}: super
  vertex algebras; the triplet and doublet algebras; graded unitarity in the sense of Arabi Ardehali,
  Beem, Lemos and Rastelli, the Lee--Yang model
  and weight-one positivity of affine $\mathfrak{sl}_2$ at level $-8/5$.
  Chapter~\ref{ch:L19} proves the positivity clauses of graded unitarity of the boundary minimal models
  $M(2,2n+1)$ through filtration degree two only; the full conjecture is not claimed. Its last section
  (Section~\ref{l19:sec:belowcritical}) proves the first full instance of the predicate
  \lean{GradedUnitary} in the library: for every real $k<-2$ the \emph{universal} affine VOA
  $V^k(\mathfrak{sl}_2)$, with its normalised form, compact real structure and current-length flag, is
  graded unitary, with no hypothesis beyond $k<-2$.
\item[Part VI, Two recent theorems, formalized.] Chapters~\ref{ch:L20}--\ref{ch:L29}. Despite its
  title, the part has ten chapters and formalizes recent research of several kinds: some chapters
  prove published theorems, and others reduce further research statements to named, unproved inputs,
  and they say so.
  \begin{itemize}
  \item Chapter~\ref{ch:L20}: the Arakawa--Dai theorems for $L_k(\mathfrak{sl}_2)$ at the boundary
    admissible levels $k=-2+2/(2m+1)$, $m\ge1$, proved with only the hypothesis $1\le m$. The vacuum
    character is proved as stated. Classical freeness, the paper's Theorem~1.1, is one Lean theorem,
    \lean{arakawaDai_theorem_1_1} (Theorem~\ref{l20:thm:ad11}): the presented algebra
    $\C[e,h,f]/(eQ^m,hQ^m,fQ^m)$ is isomorphic to the $C_2$-algebra $R_V$, and its arc algebra maps
    isomorphically onto Li's associated graded. Since the library has no functor $J_\infty$, the arc
    algebra $J_\infty R_V$ is realized as the jet algebra of the presented algebra, with its universal
    property. Its Virasoro
    companion (Section~\ref{l20:sec:vir}) proves, for every $m\ge0$ and with no hypothesis, that the
    minimal model $L(c_{2,2m+3},0)$ is classically free, with $R_L\cong\C[T]/(T^{m+1})$. A further
    section (Section~\ref{l20:sec:p3}) describes finite commutative-algebra files that prove no theorem
    about a VOA and are kept because Chapter~\ref{ch:L24} reuses them.
  \item Chapter~\ref{ch:L21}: Huang's theory. The ordinary modules of every strongly rational VOA form a
    braided monoidal category, and the logarithmic modules of a $C_2$-cofinite VOA form a braided tensor
    category. Rigidity, the Verlinde formula and the modular category are proved only for \emph{nice}
    VOAs, which the library defines as strongly rational \emph{and pseudo-unitary}; Huang's theorems need
    no pseudo-unitarity. A first step of its removal is proved: $V'\cong V$ and the vacuum-fusion chain hold for every
    strongly rational VOA (\lean{VOA.contra_vacuum_equiv_of_stronglyRational}). Rigidity is two-sided (\lean{VOA.Huang.HuangRigid.rigid_ordinary} gives
    Mathlib's \lean{RigidCategory}). The congruence property has named, hypothesis-free theorems
    (\lean{dlnCongruenceInput_holds}, \lean{zhuCongruenceInput_holds}, Theorem~\ref{l21:thm:congnamed});
    of Ng--Schauenburg's Theorem~6.8 only the kernel half is proved.
  \item Chapter~\ref{ch:L22}: the Gannon--Rayhaun classification of nice VOAs of central charge one,
    forward and converse, with no hypothesis. ``Rational'' in \lean{Nice} is the library's
    \lean{AdmissiblyRational}; for $C_2$-cofinite VOAs Dong--Li--Mason rationality implies it
    (\lean{VOA.Classification.COne.DLMRat.admissiblyRational_of_dlmRational}), and the converse is open;
    the classification and converse are stated for carriers in \lean{Type 0}.
  \item Chapter~\ref{ch:L23}: the universal affine VOA $V^k(\mathfrak g)$ for an arbitrary
    finite-dimensional Lie algebra with an invariant form and the Casimir data (``Lie data''), at every
    non-critical level; Lie data for every finite-dimensional simple $\mathfrak g$ given a
    nondegenerate symmetric invariant form; the simple quotient $L_k(\mathfrak g)$,
    simple, with weight one spanned by the currents; and concrete simple Lie algebras in Lean:
    $\mathfrak{sl}_n$, the simply-laced $\mathfrak{sl}_3,\mathfrak{so}_8,\mathfrak e_6,\mathfrak e_7,
    \mathfrak e_8$ as weight-one spaces of lattice VOAs, and $\mathfrak g_2,\mathfrak f_4,\mathfrak e_6$ from
    kernel-checked structure constants. These results are unconditional. That the algebras are of the
    named Cartan type is not claimed (the classification of simple Lie algebras is not formalized).
  \item Chapters~\ref{ch:L24}--\ref{ch:L27}, the Deligne series: classical freeness of $L_k(\mathfrak g)$
    for the seven types $A_2,G_2,D_4,F_4,E_6,E_7,E_8$ at $k=-1-h^\vee/6$, with the explicit presentation
    of $R_V$ by the quadrics $I_2^{\mathrm{ns}}$. Chapter~\ref{ch:L24} proves, with no hypothesis, the
    reduction of classical freeness to one inequality of dimensions, the explicit presentation
    ($I_2^{\mathrm{ns}}$ lies in the kernel) for all seven types, current-jet surjectivity and the
    geometric carrier; its carrier chain, certified type by type by large kernel-checked integer
    certificates, is conditional on two fusion inputs, U5 and U6. Chapter~\ref{ch:L25} proves for all
    seven types the lower bound $\dim V_N\ge c_N$, where $c_N$ are the coefficients of the normalised
    solution of the second-order Arakawa--Kawasetsu MLDE; the argument is formal, from a vacuum quartic
    relation through Zhu's recursion, and does not use the general theorem of Chapter~\ref{ch:L16}. It
    also proves a library of formal $q$-series identities (Jacobi triple and quintuple products,
    Ramanujan's differential identities, and others). Equality $\dim V_N=c_N$ is a Lean theorem for $A_2$
    only. Chapter~\ref{ch:L26} proves U6 by construction and reduces U5 to a total-dimension bound and
    Demazure realizations; both are proved for $A_2$, and are named, unproved statements for the other six types.
    Chapter~\ref{ch:L27} states the seven final theorems. The $A_2$ theorem has no hypothesis (machine-confirmed: aggregate build green on 2026-10-28, \texttt{\#print axioms} standard, as
    Chapter~\ref{ch:L27} records). The
    other six are \textbf{conditional}. $D_4$ rests on one published input, the total-dimension bound for fusion products
    (\lean{D4_classicallyFree_of_inputT}; the Demazure realizations are proved in Lean for $D_4$). One step of the published
    proof of that bound is insufficient as written for $D_4$; the bound itself is not known to fail. $G_2$, $F_4$, $E_6$,
    $E_7$, $E_8$ rest on that bound, on the Demazure realizations, and on the star bound, a conjecture supported by
    numerical evidence.
  \end{itemize}
  Chapters~\ref{ch:L28} and~\ref{ch:L29} transfer graded unitarity through the BRST coset. Their headline is a
  proved implication with one explicit hypothesis, \lean{hA}: if the simple affine VOA $L_k(\mathfrak{sl}_2)$ at the
  boundary admissible level, with its normalised form, compact conjugation and corrected current-and-Sugawara
  filtration, is graded unitary, then so is the minimal model $M(2,2n+1)$ (Chapter~\ref{ch:L28}: Lee--Yang, $n=2$,
  from level $-4/3$; Chapter~\ref{ch:L29}: every $n\ge2$). The hypothesis \lean{hA} is open for every $n$, so graded
  unitarity of $M(2,2n+1)$ beyond filtration degree two is not claimed.
\item[Part VII, The Moonshine module.] Chapters~\ref{ch:02}--\ref{ch:14}: lattices, codes and the
  Leech lattice; Heisenberg and lattice VOAs and $V_\Leech$; the involution $\theta$ and
  $V_\Leech^+$; graded dimensions and $J$; twisted modules and the module $T$; the twisted Heisenberg
  and twisted lattice modules; the intertwiners $W$ and $\overline W$; the vertex algebra on $\Vnat$
  modulo one relation, and simplicity; the Golay frame and the untwisted triality; the existence of
  the twisted triality (two chapters); the key identity and the completion of the proof.
\item[Part VIII, Verification.] Chapter~\ref{ch:15}: the verification record, the dead ends, the
  numerical evidence and the open problems.
\end{description}

\subsection{Dependencies}

Part~I is a prerequisite for everything else, and within it the order is the logical order:
Chapter~\ref{ch:L01} (formal calculus) is used by Chapter~\ref{ch:L03} (the axioms), which uses
Chapter~\ref{ch:L02} only for the Virasoro algebra; Chapter~\ref{ch:L04} builds on
Chapter~\ref{ch:L03}. In Part~II, the Virasoro chapters need the Virasoro algebra of
Chapter~\ref{ch:L02}, the affine chapters its affine Lie algebras, and the lattice VOAs its lattice
data; the rigidity chapters build on the corresponding construction chapters. Part~III needs
Chapter~\ref{ch:L04}; the module-theoretic uses of \lean{VOA.pbw} (through induced modules and
the Frenkel--Zhu tower) lie in Chapters~\ref{ch:L12} and~\ref{ch:L13}.
Part~IV needs Part~III, and Chapter~\ref{ch:L14} for modular forms.
Part~VI builds on all of Parts~I--V. Chapter~\ref{ch:L20} uses the affine $\mathfrak{sl}_2$ chapters
(\ref{ch:L08}, \ref{ch:L09}) and Li's filtration and jets (Chapter~\ref{ch:L11}); Chapters~\ref{ch:L21}
and~\ref{ch:L22} use the modules, Zhu's algebra and modularity of Parts~III and~IV, and
Chapter~\ref{ch:L22} uses Chapter~\ref{ch:L21}. Chapter~\ref{ch:L23} generalizes the constructions of
Chapter~\ref{ch:L08} (and the affine $\mathfrak{sl}_3$ files, whose chapter is in preparation) and uses
the lattice VOAs of Chapter~\ref{ch:L05}. The Deligne
chapters are read in order, Chapter~\ref{ch:L24} to Chapter~\ref{ch:L27}; they use Chapter~\ref{ch:L23}
for the algebras, Chapter~\ref{ch:L11} for Li's filtration and generator words, Chapter~\ref{ch:L20} as
the model of the squeeze argument, and Chapter~\ref{ch:L15} for Zhu's recursion (in
Chapter~\ref{ch:L25}). Chapters~\ref{ch:L28} and~\ref{ch:L29} use the graded-unitarity
predicate of Chapter~\ref{ch:L19}, the affine and Virasoro algebras of Parts~I--II, and the jet-ring and character
results that Chapter~\ref{ch:L20} shares with them.
Part~VII needs Part~I throughout --- above all the VOA axioms and the
reconstruction theorem (Chapter~\ref{ch:L03}), and the automorphisms, simplicity and invariant
forms of Chapter~\ref{ch:L04} --- together with the modules of Chapter~\ref{ch:L12} (for the
intertwiners of Chapter~\ref{ch:09} and the Zhu-algebra arguments of Chapter~\ref{ch:12}) and the
lattice data of Chapter~\ref{ch:L02}; the modular function $J$ enters in Chapter~\ref{ch:05}. Within Part~VII the chapters
are in logical order.

\subsection{Two reading paths}\label{ch01:sec:paths}

\textbf{The Moonshine path} (to verify the construction of $\Vnat$): this chapter;
Chapter~\ref{ch:L01} (at least formal distributions, modes, binomials and locality);
Chapter~\ref{ch:L03} (the axioms, the Borcherds identity, Goddard uniqueness and reconstruction);
Chapter~\ref{ch:L04} (automorphisms, ideals and simplicity, invariant forms); then
Chapters~\ref{ch:02} to~\ref{ch:14} in order, consulting Chapter~\ref{ch:L12} for modules and Zhu's
algebra when Chapters~\ref{ch:09} and~\ref{ch:12} use them; finally Chapter~\ref{ch:15}. A referee in
a hurry can start from the end: Chapter~\ref{ch:14}, Section~\ref{ch14:sec:decode}, decodes the
headline theorem clause by clause and points back to every definition it uses.

\textbf{The general path} (to learn the library as a whole): Part~I in order; then the examples of
Chapters~\ref{ch:L05} and~\ref{ch:L06}; then representation theory (Chapter~\ref{ch:L12}) and
modularity (Chapters~\ref{ch:L14}, \ref{ch:L15}); then the remaining chapters by interest.
Within Part~VI the Deligne series is long enough to need a route: Chapter~\ref{ch:L11}
(Li's filtration, jets and Section~\ref{l11:sec:genwords}), Chapter~\ref{ch:L20} (the model case,
$\mathfrak{sl}_2$), Chapter~\ref{ch:L23}, then Chapters~\ref{ch:L24}--\ref{ch:L27} in order; a referee who
wants only the final statements and their named inputs can start from Chapter~\ref{ch:L27}. For graded
unitarity, read Chapter~\ref{ch:L19}, then Chapters~\ref{ch:L28} and~\ref{ch:L29} (the BRST coset).

\subsection{What is proved, and what is conditional, in Parts IV to VI}\label{ch01:sec:newstatus}

The chapters of Parts~IV to~VI mix unconditional theorems with conditional reductions, and their
titles do not say which is which. Table~\ref{ch01:tab:newstatus} records the status of their headline
results as the chapters themselves state it. ``Unconditional'' means: \texttt{\#print axioms}-clean (Lean's
three standard axioms) and no hypothesis beyond the evident parameters (Section~\ref{ch01:sec:uncond}).
``Conditional'' means \texttt{\#print axioms}-clean with the named hypotheses listed; the conditional
theorems are \emph{not} proofs of their conclusions.

\begingroup\small
\begin{longtable}{>{\raggedright\arraybackslash}p{0.09\textwidth}>{\raggedright\arraybackslash}p{0.47\textwidth}>{\raggedright\arraybackslash}p{0.36\textwidth}}
\caption{Status of the headline results of Parts~IV to~VI.}\label{ch01:tab:newstatus}\\
\toprule
Chapter & Headline & Status\\
\midrule\endhead
\ref{ch:L15} & Zhu's modularity theorem, canonical family, for $C_2$-cofinite $V$
  (\lean{zhuModularity_canonical_of_C2Cofinite_admissiblyRational}) & unconditional, with ``rational''
  $=$ \lean{AdmissiblyRational}; Dong--Li--Mason rationality implies it under $C_2$
  (\lean{admissiblyRational_of_dlmRational}); converse open\\
\ref{ch:L16} & $L_{-4/3}(\mathfrak{sl}_2)$ is not $C_2$-cofinite and has finite-dimensional zeroth Poisson
  homology & unconditional\\
\ref{ch:L16} & Arakawa--Kawasetsu weight grading (\lean{akWeightGrading}) & unconditional, every VOA (CFT type)\\
\ref{ch:L16} & Arakawa--Kawasetsu MLDE (Theorem~5.1), one monic MLDE for all analytic blocks & conditional on \lean{EtingofFinite}
  only; order unbounded; quasi-lisseness not formalised\\
\ref{ch:L16} & the same at $L_{-4/3}(\mathfrak{sl}_2)$, vacuum insertion & unconditional; order unspecified (AK's
  order~2 not proved)\\
\ref{ch:L19} & positivity of graded unitarity of $M(2,2n+1)$ through filtration degree two &
  unconditional; the full conjecture is open\\
\ref{ch:L19} & graded unitarity of $V^k(\mathfrak{sl}_2)$, real $k<-2$ & unconditional\\
\ref{ch:L20} & Arakawa--Dai: character, and classical freeness (Theorem~1.1,
  \lean{arakawaDai_theorem_1_1}) & unconditional
  ($1\le m$)\\
\ref{ch:L20} & classical freeness of $M(2,2m+3)$ & unconditional\\
\ref{ch:L21} & braided monoidal category of ordinary modules (strongly rational); braided tensor category
  of logarithmic modules ($C_2$-cofinite) & unconditional\\
\ref{ch:L21} & rigidity, Verlinde, modular category & for nice (strongly rational and pseudo-unitary)
  VOAs only; rigidity is two-sided (\lean{rigid_ordinary}); congruence: named theorems, only the
  kernel half of Ng--Schauenburg~6.8\\
\ref{ch:L22} & Gannon--Rayhaun classification, forward and converse & unconditional, with
  ``rational'' $=$ \lean{AdmissiblyRational}; carriers in \lean{Type 0}\\
\ref{ch:L23} & $V^k(\mathfrak g)$, $L_k(\mathfrak g)$, concrete simple Lie algebras & unconditional\\
\ref{ch:L24} & reduction of classical freeness; explicit presentation; jet surjectivity; carrier &
  unconditional; the carrier chain is conditional on U5, U6\\
\ref{ch:L25} & $\dim V_N\ge c_N$ for the seven Deligne types & unconditional; equality a theorem for
  $A_2$ only\\
\ref{ch:L26} & U6; reduction of U5 & U6 proved by construction; U5 reduced to named inputs, proved for $A_2$\\
\ref{ch:L27} & classical freeness of the seven Deligne VOAs & $A_2$: no hypothesis (aggregate build
  green, 2026-10-28); $D_4$: one published input, the total-dimension bound (not formalized; one step of its proof insufficient as written); the other
  five: that bound, the Demazure realizations and the conjectural star bound\\
\ref{ch:L28} & graded unitarity of the Lee--Yang model from that of $L_{-4/3}(\mathfrak{sl}_2)$ (BRST coset) &
  conditional on \lean{hA} (graded unitarity of $L_{-4/3}(\mathfrak{sl}_2)$), open\\
\ref{ch:L29} & the same for $M(2,2n+1)$, every $n\ge2$, from the boundary admissible $L_k(\mathfrak{sl}_2)$ &
  conditional on \lean{hA} at that level, open\\
\bottomrule
\end{longtable}
\endgroup

\subsection{Where each basic notion is defined}

Table~\ref{ch01:tab:notions} lists, for each basic notion, its Lean name and the place where it is
defined in full, with its Lean source and fidelity check.

\begingroup\small
\begin{longtable}{>{\raggedright\arraybackslash}p{0.33\textwidth}>{\raggedright\arraybackslash}p{0.37\textwidth}>{\raggedright\arraybackslash}p{0.22\textwidth}}
\caption{Where the basic notions are defined.}\label{ch01:tab:notions}\\
\toprule
Notion & Lean name & Defined in\\
\midrule\endhead
formal distribution & \lean{VOA.FormalDist}, \lean{VOA.FormalDist₂} & Def.~\ref{l01:def:formaldist}, \ref{l01:def:formaldist2}\\
modes, residue, derivative & \lean{VOA.mode}, \lean{VOA.res}, \lean{VOA.deriv} & Def.~\ref{l01:def:mode}\\
formal delta function & \lean{VOA.formalDelta}, \lean{VOA.delta₂} & Def.~\ref{l01:def:delta}\\
generalised binomials, signs & \lean{VOA.gBinom}, \lean{VOA.cbz}, \lean{VOA.negOnePow} & Def.~\ref{l01:def:gbinom}\\
commutator, locality, normal ordering & \lean{VOA.bracketDist}, \lean{VOA.Local}, \lean{VOA.normalOrdered} & Def.~\ref{l01:def:bracket}\\
locality $\iff$ delta expansion & \lean{VOA.local_iff_deltaSum} & Thm.~\ref{l01:thm:deltasum}\\
super-locality & \lean{VOA.Localε} & Def.~\ref{l01:def:bracketeps}\\
field (lower truncation) & \lean{VOA.IsField} & §\ref{l03:sec:conventions}\\
Dong's lemma & \lean{VOA.dong_nonneg}, \lean{VOA.dong_neg} & Thm.~\ref{l03:thm:dong}\\
$T$-covariance & \lean{VOA.Reconstruction.DCovariant} & Def.~\ref{l03:def:dcov}\\
pre-vertex algebra (Gannon's axioms, no translation covariance) & \lean{VOA.PreVertexAlgebra} & Def.~\ref{l03:def:VA}\\
vertex algebra (with translation covariance) & \lean{VOA.VertexAlgebra} & §\ref{l03:sec:intro}\\
$n$-th product, translation & \lean{VOA.PreVertexAlgebra.smul}, \lean{VOA.PreVertexAlgebra.T} & Def.~\ref{l03:def:smulT}\\
graded VOA (FLM/Lepowsky--Li) & \lean{VOA.GradedVertexOperatorAlgebra} & Def.~\ref{l03:def:VOA}\\
vertex operator algebra (of CFT type) & \lean{VOA.VertexOperatorAlgebra} & Def.~\ref{l03:def:VOA}\\
Borcherds (Jacobi) identity & \lean{VOA.VertexOperatorAlgebra.jacobi} & Thm.~\ref{l03:thm:borcherds}\\
commutator formula & \lean{VOA.mode_commutator} & Thm.~\ref{l03:thm:comm}\\
skew-symmetry & \lean{VOA.VertexOperatorAlgebra.skew_mode} & Thm.~\ref{l03:thm:skew}\\
Goddard uniqueness & \lean{VOA.PreVertexAlgebra.goddard_uniqueness} & Thm.~\ref{l03:thm:goddard}\\
reconstruction & \lean{VOA.Reconstruction.reconstruct} & Thm.~\ref{l03:thm:reconstruct}\\
Witt and Virasoro algebras & (Chapter~\ref{ch:L02}) & Def.~\ref{l02:def:witt}, \ref{l02:def:vir}\\
affine Lie algebra, vacuum module & (Chapter~\ref{ch:L02}) & Def.~\ref{l02:def:affine}, \ref{l02:def:vacuum}\\
the PBW principle (a theorem) & \lean{VOA.pbw}, \lean{VOA.PBWPrinciple} & §\ref{l02:sec:axiom}\\
zero modes; $V_1$ Lie algebra & \lean{VOA.VertexOperatorAlgebra.zeroMode} & Def.~\ref{l04:def:zeromode}; Thm.~\ref{l04:thm:lie}\\
automorphism & \lean{VOA.Aut} & Def.~\ref{l04:def:aut}\\
ideal, simple VOA & \lean{IsIdeal}, \lean{IsSimple} & Def.~\ref{l04:def:simple}\\
quotient VOA & \lean{VOA.QuotVOA.quotientVOA} & Def.~\ref{l04:def:quot}\\
VOA homomorphism & \lean{VOA.IsVOAHom} & Def.~\ref{l04:def:hom}\\
sub-VOA, fixed-point sub-VOA & \lean{VOA.VOASubalgebra}, \lean{VOA.orbifoldVOA} & Def.~\ref{ch04:def:orbifold}\\
generated sub-VOA & \lean{VOASubalgebra.generated} & Def.~\ref{l04:def:generated}\\
invariant bilinear form & \lean{InvariantForm}, \lean{IsFullyInvariant} & Def.~\ref{l04:def:invform}\\
canonical invariant form (Li) & \lean{VOA.canonicalInvariantForm} & Thm.~\ref{l04:thm:canonical}\\
simple $\iff$ nondegenerate & \lean{VOA.LiRadical.}\allowbreak\lean{isSimple_iff_}\allowbreak\lean{leftRadical_eq_bot} & Thm.~\ref{l04:thm:lirad}\\
weak and ordinary modules & \lean{VOA.WeakModule}, \lean{VOA.OrdVModule} & Chapter~\ref{ch:L12}, §\ref{l12:sec:recall}; Def.~\ref{ch09:def:module}\\
admissible modules & \lean{VOA.AdmissibleModule} & Def.~\ref{l12:def:admissible}\\
lowest weight & \lean{VOA.IsLowestWeight} & Def.~\ref{l12:def:lowest}\\
Zhu's algebra & \lean{VOA.zhuMul}, \lean{VOA.O}, \lean{VOA.A} & §\ref{l12:sec:zhu}; Def.~\ref{ch12:def:zhumul}\\
intertwiners, fusion coefficient & \lean{VOA.Intertwiner}, \lean{VOA.fusionCoeff} & §\ref{l12:sec:fusion}; Def.~\ref{ch09:def:intertwiner}\\
twisted module & \lean{VOA.TwistedModule} & Def.~\ref{ch06:def:tm}\\
super vertex algebra & \lean{VOA.SuperVertexAlgebra} & §\ref{ch01:sec:super}; Chapter~\ref{ch:L17}\\
lattices, Leech lattice & (Chapter~\ref{ch:02}) & Def.~\ref{ch02:def:lattice}\\
$q\,J(q)$ as a power series & \lean{VOA.JShift} & Def.~\ref{ch05:def:jshift}\\
the carrier $\Vnat$, its vacuum and $\omega$ & \lean{Vnat}, \lean{vacNat}, \lean{omegaNat} & Def.~\ref{ch10:def:vnat}\\
the block vertex map & \lean{Ynat} & Def.~\ref{ch10:def:Ynat}\\
\bottomrule
\end{longtable}
\endgroup

%======================================================================
\section{A short physical orientation}\label{ch01:sec:physics}
%======================================================================

Vertex algebras were invented twice: by mathematicians (Borcherds, 1986; Frenkel, Lepowsky and
Meurman, 1988), in order to construct the Monster module $\Vnat$, and --- in a less axiomatic form
--- by physicists (Belavin, Polyakov and Zamolodchikov, 1984), as the algebraic skeleton of
two-dimensional conformal field theory (CFT). The axioms look strange until one knows where they
come from. Nothing in this section is used logically; it is here so that the definitions of Part~I
look inevitable rather than arbitrary. Each chapter's physics boxes refine the picture.

\textbf{Chiral fields.} A two-dimensional quantum field theory on the Euclidean plane
$\mathbb R^2\cong\C$ assigns to each point an operator-valued field $\phi(z)$. In a conformal field
theory the symmetries include all local holomorphic coordinate changes, and the theory splits, to a
first approximation, into a holomorphic (\emph{chiral}) half whose fields depend only on $z$ and an
antiholomorphic half. A vertex algebra is the algebraic structure of the chiral half. In
\emph{radial quantisation} the radius $|z|$ plays the role of time and the space of states lives on
a circle around the origin; a chiral field is then a Laurent series $Y(u,z) = \sum u_{(n)}z^{-n-1}$
in operators, which is why the library's basic object is a formal series with coefficients in
$\End(V)$.

\textbf{States and fields.} Every state $u$ is created by a unique field from the vacuum,
$u = \lim_{z\to0}Y(u,z)\chonevac$; this \emph{state--field correspondence} is the axiom
\lean{creation}. The vacuum's own field is the identity.

\textbf{Operator product expansion and locality.} Fields at different points commute in a local
theory. In radial ordering this becomes the statement that $Y(u,z)Y(v,w)$ and $Y(v,w)Y(u,z)$ are
expansions in two regions of one meromorphic function with a pole of finite order on the diagonal;
the coefficients of the poles are the operator product expansion (OPE)
$Y(u,z)Y(v,w)\sim\sum_{j\ge0}Y(u_{(j)}v,w)/(z-w)^{j+1}$. Algebraically, without analysis, this is
expressed by \emph{locality}: $(z-w)^N[Y(u,z),Y(v,w)] = 0$ for some $N$. The products $u_{(n)}v$,
one for each $n\in\Z$, are the OPE coefficients; a vertex algebra is a vector space with infinitely
many bilinear products subject to the identities forced by the existence of the underlying
meromorphic functions.

\textbf{Conformal symmetry.} The stress tensor $T(z) = \sum L_nz^{-n-2}$, the field of the conformal
vector $\omega$, has modes satisfying the Virasoro algebra with central charge $c$. $L_0$ measures
the conformal weight (energy) of a state, $L_{-1}$ generates translations and acts as $\partial_z$ on
fields. A vertex \emph{operator} algebra is a vertex algebra with such a stress tensor.

\textbf{The free boson and $\Vnat$.} The simplest example is the chiral free boson: the Fock space
$\C[a_{-1},a_{-2},\dots]\chonevac$ with $[a_m,a_n] = m\,\delta_{m+n,0}$ and OPE
$a(z)a(w)\sim(z-w)^{-2}$; it has $c=1$. Twenty-four free bosons compactified on the Leech torus
$\mathbb R^{24}/\Leech$ give the lattice VOA $V_\Leech$ (Chapter~\ref{ch:03}). The reflection
$x\mapsto -x$ of the torus is a symmetry; gauging it --- the $\Z/2$ \emph{orbifold} --- keeps the
reflection-invariant states of $V_\Leech$ and adds a \emph{twisted sector} of strings closed only up
to the reflection, whose bosons have half-integer modes. Dixon, Ginsparg and Harvey recognised that
the result is the Moonshine module $\Vnat$ of Frenkel, Lepowsky and Meurman: a \emph{holomorphic}
CFT of central charge $24$, with partition function $J(q) = q^{-1} + 196884\,q + \cdots$, no
weight-one states (the Leech lattice has no roots, and the orbifold removes the remaining currents),
and the Monster group as its symmetry group. The Lean construction (Chapters~\ref{ch:02}
to~\ref{ch:14}) follows this physical recipe step by step; its hardest part is to prove that the
fields of the twisted sector are local with each other, the step physicists take for granted.

\begin{physics}
\textbf{Summary of the dictionary.} States $\leftrightarrow$ vectors of $V$; local operators
$\leftrightarrow$ fields $Y(u,z)$; OPE coefficients $\leftrightarrow$ products $u_{(n)}v$; crossing
symmetry of chiral correlators $\leftrightarrow$ locality and the Borcherds identity; stress tensor
$\leftrightarrow$ conformal vector $\omega$; conformal weight $\leftrightarrow$ $L_0$-grading;
internal symmetries $\leftrightarrow$ automorphisms; sectors $\leftrightarrow$ modules; chiral vertex
operators $\leftrightarrow$ intertwining operators; orbifold $\leftrightarrow$ fixed-point subalgebra
plus twisted modules; partition function $\leftrightarrow$ graded dimension. The correspondence is
heuristic: a vertex algebra axiomatises the chiral operator algebra, not a full non-chiral CFT, and
none of the physical arguments is used as a proof anywhere in the book.
\end{physics}

%======================================================================
\section{Lean glossary}\label{ch01:sec:glossary}
%======================================================================

Table~\ref{ch01:tab:glossary} collects the symbols and idioms of this chapter, with their meaning, the
place in this chapter where they are explained, and the first chapter (in the order of the book) in
which they occur. Later chapters' fidelity checks point back here instead of re-explaining them.

\begingroup\small
\begin{longtable}{>{\raggedright\arraybackslash}p{0.28\textwidth}>{\raggedright\arraybackslash}p{0.44\textwidth}>{\raggedright\arraybackslash}p{0.10\textwidth}>{\raggedright\arraybackslash}p{0.09\textwidth}}
\caption{Lean glossary.}\label{ch01:tab:glossary}\\
\toprule
Symbol / idiom & Meaning & Here & First use\\
\midrule\endhead
\lean{x : X} & $x$ is a term of type $X$ & §\ref{ch01:sec:lean} & L01\\
\lean{Prop}, \lean{h : P} & propositions; $h$ is a proof of $P$ & §\ref{ch01:sec:lean} & L01\\
\lean{Type u}, \lean{Type*}, \lean{Type _} & a type (``set''), with universe level named, fresh or inferred & §\ref{ch01:sec:lean} & L01\\
\lean{def}, \lean{theorem}, \lean{lemma} & definition; assertion with proof & §\ref{ch01:sec:lean} & L01\\
\lean{abbrev} & definition unfolded eagerly (same logic as \lean{def}) & §\ref{ch01:sec:synonyms} & L01\\
\lean{noncomputable} & uses choice; cannot be evaluated; no logical cost & §\ref{ch01:sec:lean} & L01\\
\lean{structure}, \lean{extends}, fields, \lean{x.field} & bundle of data and proofs; inheritance; projection & §\ref{ch01:sec:structures} & L01\\
\lean{class}, \lean{instance}, \lean{[C X]} & structures found by automatic search; the library declares one, \lean{FusionCategory} & §\ref{ch01:sec:lean} & L01\\
\lean{(x : X)}, \lean{{x : X}} & explicit, implicit argument; both universally quantified & §\ref{ch01:sec:lean} & L01\\
\lean{variable} & arguments shared by the declarations of a file & §\ref{ch01:sec:lean} & L01\\
\lean{maxSynthPendingDepth} & depth of instance search; set to 3 by the build; audit trap & §\ref{ch01:sec:lean}, §\ref{ch01:sec:trust} & L01\\
\lean{rfl}, \lean{Iff.rfl}, \lean{funext} & proof by unfolding definitions; same for \lean{↔}; function extensionality & §\ref{ch01:sec:synonyms} & L01\\
\lean{@[simp]}, \lean{@[ext]}, \lean{@[inherit_doc]}, \lean{@[reducible]} & attributes: metadata for automation, no logical content; \lean{ext} applies extensionality lemmas & §\ref{ch01:sec:synonyms} & L01\\
\lean{inferInstanceAs (C X)} & the instance of \lean{C X} found by search, transported to a synonym & §\ref{ch01:sec:synonyms} & L01\\
\lean{⦃x : X⦄} & strict implicit argument; read ``for all $x$'' & §\ref{ch01:sec:lean} & L05\\
\lean{inductive}, \lean{deriving}, \lean{open X in}, \lean{scoped}, \lean{_root_.} & other declaration syntax & §\ref{ch01:sec:lean} & L01\\
\lean{where}, \lean{carrier := {x | ...}} & building a structure field by field; set-builder notation & §\ref{ch01:sec:synonyms} & L01\\
\lean{letI}, \lean{haveI} & install a structure (or fact) locally inside a definition & §\ref{ch01:sec:algebra} & L03\\
\lean{fun x => e}, \lean{⟨a, b⟩} & the function $x\mapsto e$; anonymous constructor (tuple) & §\ref{ch01:sec:lean} & L01\\
\lean{if c then a else b} & case distinction (needs \lean{Decidable c}) & §\ref{ch01:sec:numbers} & L01\\
\lean{∀}, \lean{∃}, \lean{∧}, \lean{∨}, \lean{¬}, \lean{→}, \lean{↔} & logic; order of quantifiers matters; \lean{→} associates right & §\ref{ch01:sec:lean} & L01\\
\lean{∃ a ∈ P, ...} & $\exists a,\ a\in P\wedge\dots$ & §\ref{ch01:sec:linalg} & L01\\
\lean{Nonempty X}, \lean{IsEmpty X} & $X$ has an element; $X$ has none & §\ref{ch01:sec:choice} & L02\\
\lean{Function.Injective f} & $f$ is injective & §\ref{ch01:sec:choice} & L02\\
\lean{Function.support f}, \lean{S.Finite}, \lean{S.Infinite}, \lean{Finite X} & $\{f\ne0\}$; finiteness of sets and types & §\ref{ch01:sec:sums} & L01\\
\lean{Nat.card}, \lean{Set.ncard}, \lean{s.card} & cardinalities; the first two are $0$ for infinite sets & §\ref{ch01:sec:sums} & L12\\
\lean{Set X}, \lean{Set.univ}, \lean{Set.range f}, \lean{∈}, \lean{⊆}, \lean{∩}, \lean{∪} & subsets and their operations & §\ref{ch01:sec:sums} & L02\\
\lean{List X}, \lean{Finset.filter}, \lean{Finset.antidiagonal} & finite sequences; filtered finite sets; pairs with given sum & §\ref{ch01:sec:sums} & L02\\
\lean{Module.End.eigenspace T μ}, \lean{T.maxGenEigenspace μ} & $\ker(T-\mu)$; generalised eigenspace & §\ref{ch01:sec:linalg} & L02\\
\lean{ExactAt f g} (library) & $\ker g = \operatorname{im} f$ for module maps & §\ref{ch01:sec:linalg} & L12\\
\lean{IsCompl P Q}, \lean{Nontrivial X} & $X = P\oplus Q$; $X$ has two distinct elements & §\ref{ch01:sec:linalg} & L02\\
\lean{LinearIndependent}, \lean{Module.Basis}, \lean{b.repr}, \lean{Module.Finite} & independence; bases and coordinates; finite generation & §\ref{ch01:sec:linalg} & L02\\
\lean{Submodule.Quotient.mk}, \lean{liftQ}, \lean{mkQ} & class in a quotient; induced map; quotient map & §\ref{ch01:sec:linalg} & L02\\
\lean{LinearEquiv.of...}, \lean{LinearMap.fst}, \lean{smulRight}, \lean{compl₁₂}, \lean{IsSymm}, \lean{Nondegenerate} & standard constructions of maps and forms & §\ref{ch01:sec:linalg} & L02\\
\lean{TensorProduct.mk}, \lean{.map}, \lean{.lift}, \lean{lTensor}, \lean{rTensor} & tensor-product constructions & §\ref{ch01:sec:linalg} & L02\\
\lean{DirectSum.decompose}, \lean{DirectSum.toModule}, \lean{chooseDecomposition} & decomposition isomorphism; map out of a sum; chosen decomposition & §\ref{ch01:sec:linalg} & L04\\
\lean{Finsupp.single}, \lean{Finsupp.linearCombination}, \lean{Finsupp.weight} & finitely supported constructions & §\ref{ch01:sec:linalg} & L02\\
\lean{IsUnit}, \lean{ˣ}, \lean{⁻¹}, \lean{Commute}, \lean{algebraMap}, \lean{→ₐ[ℂ]}, \lean{→+} & units, inverses, commuting elements, algebra and additive maps & §\ref{ch01:sec:algebra} & L01\\
\lean{Odd}, \lean{∣}, \lean{Nat.Coprime}, \lean{n.factorial}, \lean{n.choose k} & integer notions & §\ref{ch01:sec:algebra} & L01\\
\lean{Complex.I}, \lean{Complex.exp}, \lean{Real.pi}, \lean{Real.sqrt}, \lean{CFC.sqrt} & $i$, $\exp$, $\pi$, square roots (junk $0$ for negatives) & §\ref{ch01:sec:algebra} & L02\\
\lean{∘}, \lean{∃!}, \lean{.val}, \lean{Subtype.ext} & composition; unique existence; subtypes & §\ref{ch01:sec:algebra} & L01\\
\lean{Polynomial.X}, \lean{C}, \lean{eval}, \lean{natDegree}, \lean{derivative}; \lean{MvPolynomial.pderiv} & polynomial vocabulary (degree $0$ for $P=0$) & §\ref{ch01:sec:algebra} & L05\\
\lean{ModularForm}, \lean{qExpansion}, \lean{HasProd} & modular forms; $q$-expansions; convergent products & §\ref{ch01:sec:algebra}, §\ref{ch01:sec:sums} & ch02\\
\lean{LieRing}, \lean{LieAlgebra}, \lean{→ₗ⁅ℂ⁆}, \lean{UniversalEnvelopingAlgebra}, \lean{twoCocycle} & Lie algebras, their maps, $U(L)$, cocycles & §\ref{ch01:sec:algebra} & L02\\
\lean{Ideal}, \lean{Ideal.span}, \lean{RingHom.ker}, \lean{IsSimpleModule}, \lean{IsSemisimpleRing} & ring and module theory & §\ref{ch01:sec:algebra} & L02\\
\lean{h.choose}, \lean{Classical.choose} & a witness chosen by the axiom of choice & §\ref{ch01:sec:choice} & L01\\
\lean{Module.Free.chooseBasis} & an arbitrary basis of a free module & §\ref{ch01:sec:choice} & L05\\
\lean{ℕ}, \lean{ℤ}, \lean{ℂ}, \lean{(n : ℂ)}, \lean{↑n} & number types ($0\in\N$); casts & §\ref{ch01:sec:numbers} & L01\\
\lean{m - n} on \lean{ℕ}, \lean{x / 0}, \lean{k.toNat} & junk values: truncated, $0$, $\max(k,0)$ & §\ref{ch01:sec:junk} & L01\\
\lean{ZMod n}, \lean{Fin n} & $\Z/n\Z$; $\{0,\dots,n-1\}$ & §\ref{ch01:sec:numbers} & L01\\
\lean{Even n}, \lean{negOnePow n} & $n$ even; $(-1)^n$ & §\ref{ch01:sec:numbers} & L01\\
\lean{FormalDist 𝒲}, \lean{𝒲⟦z^±⟧} & opaque synonym of \lean{ℤ → 𝒲}: formal series & §\ref{ch01:sec:synonyms} & L01\\
\lean{mode f n}, \lean{mode (Y u) n} & $f(-n-1)$; the operator $u_{(n)}$ & Ex.~\ref{ch01:ex:mode} & L01\\
\lean{1} on \lean{FormalDist} & the identity field & Ex.~\ref{ch01:ex:instOne} & L01\\
\lean{odd k ↦ k/2} encoding & half-integer modes indexed by odd integers; even $k$ junk & §\ref{ch01:sec:junk} & ch06\\
\lean{→ₗ[ℂ]}, \lean{≃ₗ[ℂ]}, \lean{LinearMap.id} & linear maps; linear isomorphisms; identity & §\ref{ch01:sec:linalg} & L01\\
\lean{End ℂ 𝒱}, \lean{*}, \lean{⁅a, b⁆}, \lean{D ^ N} & endomorphisms; composition; commutator; power & §\ref{ch01:sec:linalg} & L01\\
\lean{M →ₗ[ℂ] M →ₗ[ℂ] ℂ}, \lean{LinearMap.BilinForm} & bilinear (not sesquilinear) forms & §\ref{ch01:sec:linalg} & L02\\
\lean{Module.Dual ℂ X} & dual space & §\ref{ch01:sec:linalg} & L06\\
\lean{Basis ι ℂ X} & a basis indexed by $\iota$ & §\ref{ch01:sec:linalg} & L02\\
\lean{Submodule ℂ 𝒱}, \lean{≤}, \lean{⊥}, \lean{⊤} & subspaces; inclusion; $0$; whole space & §\ref{ch01:sec:linalg} & L01\\
\lean{⊓}, \lean{⊔}, \lean{⨆ i, S i} & intersection; sum; sum of a family & §\ref{ch01:sec:linalg} & L04\\
\lean{Submodule.span}, \lean{ℂ ∙ v} & span; the line $\C v$ & §\ref{ch01:sec:linalg} & L02\\
\lean{LinearMap.ker}, \lean{LinearMap.range}, \lean{comap}, \lean{map}, \lean{subtype} & kernel; image; preimage; image; inclusion & §\ref{ch01:sec:linalg} & L04\\
\lean{↥S} & a subspace as a vector space (coercion) & §\ref{ch01:sec:linalg} & L04\\
\lean{X ⧸ S} & quotient space & §\ref{ch01:sec:linalg} & L02\\
\lean{ℤ → Submodule ℂ 𝒱}, \lean{DirectSum.IsInternal} & grading; $V = \bigoplus_n V_n$ & §\ref{ch01:sec:linalg} & L02\\
\lean{DirectSum ℂ X} & external direct sum & §\ref{ch01:sec:linalg} & L12\\
\lean{Module.finrank}, \lean{FiniteDimensional} & dimension ($0$ if infinite); finite-dimensionality & §\ref{ch01:sec:linalg} & L02\\
\lean{A × B}, \lean{(a, b)}, \lean{x.1}, \lean{S.prod T} & product space; pair; projection; product subspace & §\ref{ch01:sec:linalg} & L01\\
\lean{A ⊗[ℂ] B}, \lean{a ⊗ₜ[ℂ] b} & tensor product; pure tensor & §\ref{ch01:sec:linalg} & L02\\
\lean{Matrix ι ι ℂ}, \lean{M.det} & square matrix; determinant & §\ref{ch01:sec:linalg} & L01\\
\lean{∑ i ∈ s, f i}, \lean{Finset.range N}, \lean{Finset.univ} & finite sums; $\{0,\dots,N-1\}$; all elements & §\ref{ch01:sec:sums} & L01\\
\lean{∑ᶠ i, f i} (\lean{finsum}) & sum with finite support; $0$ if support infinite & §\ref{ch01:sec:sums} & L01\\
\lean{ι →₀ R}, \lean{μ.sum g} & finitely supported functions; sum over support & §\ref{ch01:sec:sums} & L02\\
\lean{∑'} (\lean{tsum}) & infinite sum; $0$ if not summable & §\ref{ch01:sec:sums} & L14\\
\lean{Polynomial ℂ}, \lean{MvPolynomial σ ℂ}, \lean{aeval} & polynomials; several variables; evaluation & §\ref{ch01:sec:algebra} & L05\\
\lean{PowerSeries ℂ}, \lean{PowerSeries.coeff n} & $\C[[q]]$; coefficient of $q^n$ & §\ref{ch01:sec:algebra} & L05\\
\lean{Subgroup G}, \lean{Subgroup.zpowers g} & subgroups; cyclic subgroup $\langle g\rangle$ & §\ref{ch01:sec:algebra} & ch04\\
\lean{sorry}, \lean{sorryAx}, \lean{axiom} & unproved goal (forbidden); postulate & §\ref{ch01:sec:trust} & L01\\
\texttt{\#print axioms} & lists the axioms a proof uses (not its hypotheses) & §\ref{ch01:sec:trust} & L01\\
\bottomrule
\end{longtable}
\endgroup

%======================================================================
\section{Referee's checklist}\label{ch01:sec:checklist}
%======================================================================

This chapter quotes only a few declarations, as reading examples; their full fidelity checks are in
the chapters indicated.

\begingroup\small
\begin{longtable}{>{\raggedright\arraybackslash}p{0.25\textwidth}>{\raggedright\arraybackslash}p{0.33\textwidth}>{\raggedright\arraybackslash}p{0.2\textwidth}>{\raggedright\arraybackslash}p{0.12\textwidth}}
\toprule
Lean declaration & Intended meaning & Hypotheses / status & Check\\
\midrule\endhead
\lean{VOA.mode} & $f_{(m)} = f(-m-1)$ & definition; no axioms & Ex.~\ref{ch01:ex:mode}; Def.~\ref{l01:def:mode}\\
\lean{VOA.FormalDist.instOne} & \lean{1} is the identity field & definition; no axioms & Ex.~\ref{ch01:ex:instOne}\\
\lean{VOA.Aut} & automorphisms of a VOA & definition & Ex.~\ref{ch01:ex:aut}; Def.~\ref{l04:def:aut}\\
\lean{VOA.VertexOperatorAlgebra}, \lean{…VertexOperatorAlgebra.L} & VOA of CFT type; $L_n = \omega_{(n+1)}$ & definition (CFT type by definition; the FLM notion is \lean{GradedVertexOperatorAlgebra}) & Ex.~\ref{ch01:ex:VOA}; Def.~\ref{l03:def:VOA}\\
\lean{VOA.SuperVertexAlgebra} & super vertex algebra & definition & §\ref{ch01:sec:super}; Chapter~\ref{ch:L17}\\
\lean{VOA.Moonshine.moonshine} & $\Vnat$ exists with the FLM data, $c=24$, $V_1=0$, $\dim V_2=196884$, character $J$, simple & no hypotheses; unconditional; standard axioms & §\ref{ch01:sec:headline}; §\ref{ch14:sec:decode}\\
\lean{VOA.pbw} & two PBW consequences & theorem, unconditional (formerly a project axiom); not used by Moonshine & §\ref{ch01:sec:trust}; §\ref{l02:sec:axiom}\\
\bottomrule
\end{longtable}
\endgroup

%======================================================================
\section{Notes and references}\label{ch01:sec:notes}
%======================================================================

Lean~4 is described in L.~de~Moura and S.~Ullrich, \emph{The Lean~4 theorem prover and programming
language}, in: Automated Deduction --- CADE~28, Lecture Notes in Computer Science \textbf{12699},
Springer, 2021, 625--635; Mathlib in The mathlib Community, \emph{The Lean mathematical library},
Proceedings of the 9th ACM SIGPLAN International Conference on Certified Programs and Proofs (CPP
2020), 367--381 (arXiv:1910.09336). A gentle introduction for mathematicians is J.~Avigad and
P.~Massot, \emph{Mathematics in Lean}, available with the Lean community documentation.

The definition of a vertex algebra is due to R.~E.~Borcherds, \emph{Vertex algebras, Kac--Moody
algebras, and the Monster}, Proc.\ Nat.\ Acad.\ Sci.\ USA \textbf{83} (1986) 3068--3071. The
definition of a vertex operator algebra and the construction of the Moonshine module are in
I.~B.~Frenkel, J.~Lepowsky and A.~Meurman, \emph{Vertex Operator Algebras and the Monster}, Pure and
Applied Mathematics \textbf{134}, Academic Press, 1988. Standard textbooks are V.~Kac, \emph{Vertex
Algebras for Beginners}, University Lecture Series \textbf{10}, AMS, 2nd ed.\ 1998; E.~Frenkel and
D.~Ben-Zvi, \emph{Vertex Algebras and Algebraic Curves}, Mathematical Surveys and Monographs
\textbf{88}, AMS, 2nd ed.\ 2004; and J.~Lepowsky and H.~Li, \emph{Introduction to Vertex Operator
Algebras and Their Representations}, Progress in Mathematics \textbf{227}, Birkh\"auser, 2004. The
library's files follow, and their docstrings cite section by section, T.~Gannon, \emph{Moonshine
beyond the Monster}, Cambridge University Press, 2006, Chapter~5.

On the physics side, the founding paper of two-dimensional conformal field theory is A.~A.~Belavin,
A.~M.~Polyakov and A.~B.~Zamolodchikov, \emph{Infinite conformal symmetry in two-dimensional quantum
field theory}, Nucl.\ Phys.\ B \textbf{241} (1984) 333--380; a textbook account is P.~Di~Francesco,
P.~Mathieu and D.~S\'en\'echal, \emph{Conformal Field Theory}, Springer, 1997. The orbifold
interpretation of $\Vnat$ is due to L.~Dixon, P.~Ginsparg and J.~Harvey, \emph{Beauty and the beast:
superconformal symmetry in a Monster module}, Commun.\ Math.\ Phys.\ \textbf{119} (1988) 221--241.

The project-internal sources for this chapter are the files cited in the text, the
repository's \texttt{README.md} (the policy on forbidden constructs, and the
former single project axiom, now the theorem \lean{VOA.pbw}), the book's governing document \texttt{Notes/Textbook/CONSTITUTION.md}, the scripts
\texttt{scripts/textbook/sync.py} and \texttt{scripts/textbook/cite\_paths.py}, and the catalogue of defects in Chapter~\ref{ch:15}. The
release-audit reports \texttt{Notes/Textbook/referee/RELEASE\_AUDIT\_L16.md}, \dots, \texttt{RELEASE\_AUDIT\_L22.md} are
the source of the corrected statements of Chapters~\ref{ch:L16}, \ref{ch:L20}, \ref{ch:L21} and~\ref{ch:L22}
recorded in Section~\ref{ch01:sec:map}. The
definitions that an earlier version of this chapter presented in full --- formal calculus, the
vertex algebra axioms, the Borcherds identity, reconstruction, modules, automorphisms, ideals,
invariant forms, Zhu's algebra and intertwiners --- are now presented in Chapters~\ref{ch:L01},
\ref{ch:L03}, \ref{ch:L04} and~\ref{ch:L12}, and the orbifold subalgebra in Chapter~\ref{ch:04}.

\endgroup
